%  LaTeX support: latex@mdpi.com 
%  For support, please attach all files needed for compiling as well as the log file, and specify your operating system, LaTeX version, and LaTeX editor.
% Primary target: MDPI Mathematics.  The file also has a standalone fallback so
% it can be compiled before the MDPI Definitions/ folder is supplied.
%=================================================================
%\documentclass[journal,article,submit,moreauthors]
\newif\ifmdpiclass
\IfFileExists{Definitions/mdpi.cls}{%
  \mdpiclasstrue
  \documentclass[jimaging,article,submit,pdftex,moreauthors]{Definitions/mdpi}%
}{%
  \mdpiclassfalse
  \documentclass[11pt,a4paper]{article}%
}
% For posting an early version of this manuscript as a preprint, you may use "preprints" as the journal. Changing "submit" to "accept" before posting will remove line numbers.
\ifmdpiclass
\else
  \usepackage[a4paper,margin=25mm]{geometry}
  \usepackage{graphicx}
  \usepackage{booktabs}
  \usepackage{tabularx}
  \usepackage{array}
  \usepackage{float}
  \usepackage{amsmath,amssymb}
  \usepackage{enumitem}
  \usepackage[numbers,sort&compress]{natbib}
  \usepackage{url}
  % The standalone fallback avoids a hard dependency on hyperref so it can be
  % compiled on minimal TeX installations. The MDPI class supplies hyperlinks.
  \providecommand{\href}[2]{#2}
\fi
% The following journals will use the APA reference format:
% admsci, aieduc, behavsci, businesses, econometrics, economies, education, ejihpe, famsci, games, humans, ijcs, ijfs, jintelligence, journalmedia, jrfm, jsam, languages, peacestud, psycholint, publications, tourismhosp, youth
% The following journals will use the Chicago reference format:
% arts, genealogy, histories, humanities, laws, literature, religions, risks, socsci
%--------------------
% Class Options:
%--------------------
%----------
% journal
%----------
% Choose between the following MDPI journals:
% accountaudit, acoustics, actuators, addictions, adhesives, admsci, adolescents, aerobiology, aerospace, agriculture, agriengineering, agrochemicals, agronomy, ai, aichem, aieduc, aieng, aimater, aimed, aipa, air, aisens, algorithms, allergies, alloys, amh, analog, analytica, analytics, anatomia, anesthres, animals, antibiotics, antibodies, antioxidants, applbiosci, appliedchem, appliedmath, appliedphys, applmech, applmicrobiol, applnano, applsci, aquacj, architecture, arm, arthropoda, arts, asc, asi, astronautics, astronomy, atmosphere, atoms, audiolres, automation, axioms, bacteria, batteries, bdcc, behavsci, beverages, biochem, bioengineering, biologics, biology, biomass, biomechanics, biomed, biomedicines, biomedinformatics, biomimetics, biomolecules, biophysica, bioresourbioprod, biosensors, biosphere, biotech, birds, blockchains, bloods, blsf, brainsci, breath, buildings, businesses, cancers, carbon, cardio, cardiogenetics, cardiovascmed, catalysts, cells, ceramics, challenges, chemengineering, chemistry, chemosensors, chemproc, children, chips, cimb, civileng, cleantechnol, climate, clinbioenerg, clinpract, clockssleep, cmd, cmtr, coasts, coatings, colloids, colorants, commodities, complexities, complications, compounds, computation, computers, condensedmatter, conservation, constrmater, cosmetics, covid, crops, cryo, cryptography, crystals, csmf, ctn, culture, curroncol, cyber, dairy, data, ddc, dentistry, dermato, dermatopathology, designs, devices, dhi, diabetology, diagnostics, dietetics, digital, disabilities, diseases, diversity, dna, drones, dynamics, earth, ebj, ecm, ecologies, econometrics, economies, edm, education, eesp, ejihpe, electricity, electrochem, electronicmat, electronics, encyclopedia, endocrines, energies, eng, engproc, entomology, entropy, environments, environremediat, epidemiologia, epigenomes, esa, est, famsci, fermentation, fibers, fintech, fire, fishes, fluids, foods, forecasting, forensicsci, forests, fossstud, foundations, fractalfract, fuels, future, futureinternet, futurepharmacol, futurephys, futuretransp, galaxies, games, gases, gastroent, gastrointestdisord, gastronomy, gels, genealogy, genes, geographies, geohazards, geomatics, geometry, geosciences, geotechnics, geriatrics, germs, glacies, grasses, green, greenhealth, gucdd, hardware, hazardousmatters, healthcare, hearts, hemato, hematolrep, hep, heritage, higheredu, highthroughput, histories, horticulturae, hospitals, humanities, humans, hydrobiology, hydrogen, hydrology, hydropower, hygiene, idr, iic, ijcs, ijem, ijerph, ijfs, ijgi, ijmd, ijms, ijns, ijom, ijpb, ijt, ijtm, ijtpp, ime, immuno, informatics, information, infrastructures, inorganics, insects, instruments, inventions, iot, j, jaestheticmed, jal, jcdd, jcm, jcp, jcrm, jcs, jcto, jdad, jdb, jdream, jemr, jeta, jfb, jfmk, jgbg, jgg, jimaging, jintelligence, jlpea, jmahp, jmmp, jmms, jmp, jmse, jne, jnt, jof, joi, joitmc, joma, jop, joptm, jor, journalmedia, jox, jpbi, jphytomed, jpm, jrfm, jsam, jsan, jtaer, jvd, jzbg, kidneydial, kinasesphosphatases, knowledge, labmed, laboratories, lae, land, languages, laws, life, lights, limnolrev, lipidology, liquids, literature, livers, logics, logistics, lubricants, lymphatics, machines, macromol, magnetism, magnetochemistry, make, marinedrugs, materials, materproc, mathematics, mca, measurements, medicina, medicines, medsci, membranes, merits, metabolites, metals, meteorology, methane, metrics, metrology, micro, microarrays, microbiolres, microelectronics, micromachines, microorganisms, microplastics, microwave, minerals, mining, mmphys, modelling, molbank, molecules, mps, msf, mti, multimedia, muscles, nanoenergyadv, nanomanufacturing, nanomaterials, ncrna, ndt, network, neuroglia, neuroimaging, neurolint, neurosci, nitrogen, notspecified, nursrep, nutraceuticals, nutrients, obesities, occuphealth, oceans, ohbm, onco, optics, oral, organics, organoids, osteology, oxygen, pandemics, parasites, parasitologia, particles, pathogens, pathophysiology, peacestud, pediatrrep, pets, pharmaceuticals, pharmaceutics, pharmacoepidemiology, pharmacy, philosophies, photochem, photonics, phycology, physchem, physics, physiologia, plants, plasma, platforms, pollutants, polymers, polysaccharides, populations, poultry, powders, precipitation, precisoncol, preprints, proceedings, processes, prosthesis, proteomes, psf, psychiatryint, psychoactives, psycholint, publications, purification, quantumrep, quaternary, qubs, radiation, rdt, reactions, realestate, receptors, recycling, regeneration, religions, remotesensing, reports, reprodmed, resources, rheumato, risks, rjpm, robotics, rsee, ruminants, safety, sci, scipharm, sclerosis, seeds, sensors, separations, sexes, shi, signals, sinusitis, siuj, skins, smartcities, sna, societies, socsci, software, soilsystems, solar, solids, spectroscj, sports, standards, stats, std, stratsediment, stresses, surfaces, surgeries, suschem, sustainability, symmetry, synbio, systems, tae, targets, taxonomy, technologies, telecom, test, textiles, thalassrep, therapeutics, thermo, timespace, tomography, tourismhosp, toxics, toxins, tph, transplantology, transportation, traumacare, traumas, tri, tropicalmed, universe, urbansci, uro, vaccines, vehicles, venereology, vetsci, vibration, virtualworlds, viruses, vision, waste, water, welding, wem, wevj, wild, wind, women, world, youth, zoonoticdis
%---------
% article
%---------
% The default type of manuscript is "article", but can be replaced by: 
% abstract, addendum, article, benchmark, book, bookreview, briefcommunication, briefreport, casereport, changes, clinicopathologicalchallenge, comment, commentary, communication, conceptpaper, conferenceproceedings, correction, conferencereport, creative, datadescriptor, discussion, entry, expressionofconcern, extendedabstract, editorial, essay, erratum, fieldguide, hypothesis, interestingimages, letter, meetingreport, monograph, newbookreceived, obituary, opinion, proceedingpaper, projectreport, reply, retraction, review, perspective, protocol, shortnote, studyprotocol, supfile, systematicreview, technicalnote, viewpoint, guidelines, registeredreport, tutorial,  giantsinurology, urologyaroundtheworld
% supfile = supplementary materials
%----------
% submit
%----------
% The class option "submit" will be changed to "accept" by the Editorial Office when the paper is accepted. This will only make changes to the front page (e.g., the journal's logo will become visible), the headings, and the copyright information. Also, line numbering will be removed. Journal info and pagination for accepted papers will also be assigned by the Editorial Office.
%------------------
% moreauthors
%------------------
% If there is only one author,r the class option oneauthor should be used. Otherwise, use the class option moreauthors.
%---------
% pdftex
%---------
% The option pdftex is for use with pdfLaTeX. Remove "pdftex" for (1) compiling with LaTeX & dvi2pdf (if eps figures are used) or for (2) compiling with XeLaTeX.
%=================================================================
% MDPI internal commands - do not modify
\ifmdpiclass
\firstpage{1}
\makeatletter
\setcounter{page}{\@firstpage}
\makeatother
\pubvolume{1}
\issuenum{1}
\articlenumber{0}
\pubyear{2026}
\copyrightyear{2026}
\fi
\ifmdpiclass
%\externaleditor{Firstname Lastname} % More than 1 editor, please add `` and '' before the last editor name
\datereceived{ } 
\daterevised{ } % Comment out if no revised date
\dateaccepted{ } 
\datepublished{ } 
%\datecorrected{} % For corrected papers: "Corrected: XXX" date in the original paper.
%\dateretracted{} % For retracted papers: "Retracted: XXX" date in the original paper.
%\doinum{} % Used for some special journals, like molbank
%\pdfoutput=1 % Uncommented for upload to arXiv.org
%\CorrStatement{yes}  % For updates
%\longauthorlist{yes} % For many authors that exceed the left citation part
%\IsAssociation{yes} % For association journals
\fi
%=================================================================
% Add packages and commands here. The following packages are loaded in our class file: fontenc, inputenc, calc, indentfirst, fancyhdr, graphicx, epstopdf, lastpage, ifthen, float, amsmath, amssymb, lineno, setspace, enumitem, mathpazo, booktabs, titlesec, etoolbox, tabto, xcolor, colortbl, soul, multirow, microtype, tikz, totcount, changepage, attrib, upgreek, array, tabularx, pbox, ragged2e, tocloft, marginnote, marginfix, enotez, amsthm, natbib, hyperref, cleveref, scrextend, url, geometry, newfloat, caption, draftwatermark, seqsplit
% cleveref: load \crefname definitions after 
% Preserve supplied figures when available; flag missing assets visibly.
% \begin{document}
%=================================================================
% Please use the following mathematics environments: Theorem, Lemma, Corollary, Proposition, Characterisation, Property, Problem, Example, Examples and Definitions, Hypothesis, Remark, Definition, Notation, Assumption
%% For proofs, please use the proof environment (the amsthm package is loaded by the MDPI class).
% Standalone fallbacks for MDPI front-matter commands. These definitions are
% inactive when Definitions/mdpi.cls is available.
\ifmdpiclass
\else
  \newcommand{\Title}[1]{\gdef\papertitle{#1}}
  \newcommand{\Author}[1]{\gdef\paperauthor{#1}}
  \newcommand{\AuthorNames}[1]{}
  \newcommand{\address}[1]{\gdef\paperaddress{#1}}
  \newcommand{\corres}[1]{\gdef\papercorres{#1}}
  \renewcommand{\abstract}[1]{\gdef\paperabstract{#1}}
  \newcommand{\keyword}[1]{\gdef\paperkeywords{#1}}
\fi
%=================================================================
% Full title of the paper (Capitalised)
% Keep the title free of manual line breaks so hyperref can create valid PDF
% metadata. The MDPI class will wrap the displayed title automatically.
\Title{Refining Retinal Vessel Segmentation with Conditional Diffusion Models:\\
A Lightweight Enhancement to U-Net-Based Methods}
% Add verified ORCID identifiers with the MDPI \orcidA{}, \orcidB{}, ... commands if applicable.
\Author{Babak Bayanian $^{1}$, Hanieh Naderi $^{1,*}$ and Chinthaka Dinesh $^{2,3,*}$}
\AuthorNames{Babak Bayanian, Hanieh Naderi and Chinthaka Dinesh}
\address{%
$^{1}$ \quad College of Interdisciplinary Science and Technologies, University of Tehran, Tehran, Iran; babak.bayanian@ut.ac.ir (B.B.); hanieh.naderi@ut.ac.ir (H.N.)\\
$^{2}$ \quad Department of Mechanical and Industrial Engineering, College of Engineering, Northeastern University, Vancouver, BC, Canada; c.herathgedara@northeastern.edu (C.D.)\\
$^{3}$ \quad School of Engineering Science, Simon Fraser University, Burnaby, BC, Canada
}
\corres{Correspondence: c.herathgedara@northeastern.edu; hanieh.naderi@ut.ac.ir}
% Current address and/or shared authorship
%\firstnote{Current address: Affiliation.}  
% Current address should not be the same as any items in the Affiliation section.
%\secondnote{These authors contributed equally to this work.}
% The commands \thirdnote{} till \eighthnote{} are available for further notes.
%\simplesumm{} % Simple summary
%\conference{} % An extended version of a conference paper
% Abstract (Do not insert blank lines, i.e. \\) 
\abstract{
Retinal vessel segmentation supports quantitative vascular analysis but remains difficult in low-contrast regions and thin branches. We investigate a conditional diffusion refiner for coarse probability maps from a lightweight U-Net with reverse attention (LU-Net+RA). A frozen mask autoencoder produces a grid with fourfold spatial downsampling. A 128-dimensional global RGB embedding conditions the denoiser, and a frozen decoder reconstructs the refined map after Denoising Diffusion Implicit Model (DDIM) sampling. The refiner applies continuous Gaussian diffusion to the encoded coarse prediction rather than an encoded ground-truth target. The archived implementation contains 535{,}425 refinement parameters, including the frozen autoencoder. At 50 DDIM transitions, Dice point estimates are 0.835, 0.835, and 0.793 on the stated CHASE-DB1, DRIVE, and HRF partitions. The corresponding LU-Net+RA baselines are 0.8001, 0.7879, and 0.7325, giving absolute increases of 0.0349, 0.0471, and 0.0605. Protocol matching is author-confirmed but has not been independently audited against the original diffusion outputs; published scores provide literature context only. On CHASE-DB1, Dice increases from 0.814 to 0.836 as transitions increase from 10 to 100, illustrating an adjustable accuracy--cost trade-off. The design is parameter-efficient, while verified original paired outputs and repeated training runs are needed to assess significance and robustness.}
% Keywords
\keyword{retinal vessel segmentation; fundus imaging; diffusion refinement; spatially downsampled multi-channel representation; DDIM; reverse attention; coarse-to-fine refinement.}
% The fields PACS, MSC, and JEL may be left empty or commented out if not applicable
%\PACS{J0101}
%\MSC{}
%\JEL{}
%%%%%%%%%%%%%%%%%%%%%%%%%%%%%%%%%%%%%%%%%%
% Only for the journal Diversity
%\LSID{\url{http://}}
%%%%%%%%%%%%%%%%%%%%%%%%%%%%%%%%%%%%%%%%%%
% Only for the journal Applied Sciences
%\featuredapplication{Authors are encouraged to provide a concise description of the specific application or a potential application of the work. This section is not mandatory.}
%%%%%%%%%%%%%%%%%%%%%%%%%%%%%%%%%%%%%%%%%%
%%%%%%%%%%%%%%%%%%%%%%%%%%%%%%%%%%%%%%%%%%
% Only for the journal Data
%\dataset{DOI number or link to the deposited data set if the data set is published separately. If the data set is to be published as a supplement to this paper, this field will be filled by the journal editors. In this case, please submit the data set as a supplement.}
%\datasetlicense{License under which the data set is made available (CC0, CC-BY, CC-BY-SA, CC-BY-NC, etc.)}
%%%%%%%%%%%%%%%%%%%%%%%%%%%%%%%%%%%%%%%%%%
% Only for the journal BioTech, Fishes, Neuroimaging and Toxins
%\keycontribution{The breakthroughs or highlights of the manuscript. Authors can write one or two sentences to describe the most important part of the paper.}
%%%%%%%%%%%%%%%%%%%%%%%%%%%%%%%%%%%%%%%%%%
% Only for the journal Encyclopedia
%\encyclopediadef{For entry manuscripts only: please provide a brief overview of the entry title instead of an abstract.}
%%%%%%%%%%%%%%%%%%%%%%%%%%%%%%%%%%%%%%%%%%
% Different journals have different requirements. Please check the specific journal guidelines in the "Instructions for Authors" on the journal's official website.
%\addhighlights{yes}
%\renewcommand{\addhighlights}{%
%
%\noindent The goal is to increase the discoverability and readability of the article via search engines and other scholars. Highlights should not be a copy of the abstract, but a simple text allowing the reader to quickly and simply find out what the article is about and what can be cited from it. Each of these parts should be devoted to up to 2 bullet points.\vspace{3pt}\\
%\textbf{What are the main findings?}
% \begin{itemize}[labelsep=2.5mm,topsep=-3pt]
% \item First bullet.
% \item Second bullet.
% \end{itemize}\vspace{3pt}
%\textbf{What are the implications of the main findings?}
% \begin{itemize}[labelsep=2.5mm,topsep=-3pt]
% \item First bullet.
% \item Second bullet.
% \end{itemize}
%}
% The MDPI running head is taller than fancyhdr's default. Twenty points
% satisfies both first-page and subsequent-page header requirements; the
% top-margin adjustment preserves the original text-block position.
\ifmdpiclass
  \setlength{\headheight}{20pt}
  \addtolength{\topmargin}{-8pt}
\fi
%%%%%%%%%%%%%%%%%%%%%%%%%%%%%%%%%%%%%%%%%%
\providecommand{\dataavailability}[1]{\section*{Data Availability Statement}#1}
\providecommand{\authorcontributions}[1]{\section*{Author Contributions}#1}
\providecommand{\funding}[1]{\section*{Funding}#1}
\providecommand{\institutionalreview}[1]{\section*{Institutional Review Board Statement}#1}
\providecommand{\informedconsent}[1]{\section*{Informed Consent Statement}#1}
\providecommand{\acknowledgments}[1]{\section*{Acknowledgments}#1}
\providecommand{\conflictsofinterest}[1]{\section*{Conflicts of Interest}#1}
\newcommand{\manuscriptimage}[2][]{%
  \IfFileExists{#2}{\includegraphics[#1]{#2}}{%
    \PackageWarning{manuscript}{Missing figure: #2}%
    \fbox{\parbox[c][35mm][c]{0.90\linewidth}{\centering
      Figure file missing.\\[2mm]\texttt{\detokenize{#2}}}}%
  }%
}
\begin{document}



\ifmdpiclass

\else

\begin{center}

  {\LARGE\bfseries \papertitle\par}

  \vspace{0.8em}

  {\large \paperauthor\par}

  \vspace{0.6em}

  {\small \paperaddress\par}

  \vspace{0.3em}

  {\small \papercorres\par}

\end{center}

\vspace{0.8em}

\noindent\textbf{Abstract.} \paperabstract\par

\vspace{0.5em}

\noindent\textbf{Keywords:} \paperkeywords\par

\vspace{1em}

\fi







\section{Introduction}



Retinal vessel segmentation supports quantitative analysis of vascular morphology in fundus images and may assist research on ocular and systemic disease. Convolutional neural network (CNN)-based methods have become a common approach to this task, including architectures such as U-Net and its variants, volumetric models such as V-Net~\cite{Milletari2016}, and context-aware designs such as CE-Net~\cite{Gu2019CENet}. Despite these advances, several technical challenges remain.



A primary difficulty is that vessel pixels occupy a relatively small part of a fundus image, producing substantial class imbalance. Under this imbalance, standard pixel-wise losses can underemphasise small structures, motivating alternatives such as focal loss~\cite{Lin2017}. Thin vessels and boundaries may also have low contrast, increasing the risk of missed or over-smoothed predictions. In addition, iterative refinement introduces a practical trade-off between predictive performance and computational cost.



Lightweight architectures such as LU-Net with reverse attention (LU-Net+RA) address some of these challenges through depth-wise separable convolutions and attention to missed regions~\cite{Hernandez2025}. A single feed-forward prediction cannot revise itself. A separate refiner can therefore revisit missed vessel pixels and uncertain boundaries while retaining a fixed coarse predictor. We study a conditional diffusion refiner that receives an LU-Net+RA probability map and the corresponding fundus image, then produces a revised vessel-probability map.



We introduce a conditional diffusion refiner that operates on a spatially downsampled multi-channel representation. Denoising diffusion probabilistic models (DDPMs) support iterative generation and restoration~\cite{Ho2020,Dhariwal2021}; latent diffusion applies these updates within an autoencoder representation~\cite{Rombach2022}. Our representation contains four times as many scalar values as the input mask. Its parameter efficiency therefore concerns the model budget, rather than representation compression. We evaluate segmentation on three retinal datasets and use an explicitly identified published LU-Net+RA configuration for literature context. We also examine how the number of Denoising Diffusion Implicit Model (DDIM) transitions affects accuracy and refinement cost.



\paragraph{Contributions.}

This work makes the following contributions:



\begin{enumerate}

  %\item We propose a conditional diffusion module operating on a spatially downsampled multi-channel representation that operates on coarse LU-Net+RA probability maps, combines the encoded mask with a global red--green--blue (RGB) image embedding, and returns a refined vessel-probability map through a frozen mask decoder.



  \item We propose a conditional diffusion module operating on a spatially downsampled multi-channel representation that operates on

coarse LU-Net+RA probability maps by applying continuous Gaussian diffusion to

the encoded coarse prediction, rather than to an encoded ground-truth target;

it combines the encoded mask with a global red--green--blue (RGB) image

embedding and returns a refined vessel-probability map through a frozen mask

decoder.



  \item We report a consistent parameter audit of the released LU-Net+RA baseline and both available refinement implementations, including frozen components and conditioning layers. The archived implementation contains 535,425 refinement parameters, while the documented new reproduction contains 527,233; their complete systems contain 2,482,386 and 2,474,194 parameters, respectively.



  %\item We report segmentation point estimates on CHASE-DB1, DRIVE, and HRF and place them beside the same published LU-Net+RA ablation configuration on all three datasets. The comparison is labelled as cross-study context because the partitions and evaluation procedures are not identical.

  \item We report segmentation point estimates on CHASE-DB1, DRIVE, and HRF and

place them beside the AdamW, Dice-loss, reverse-attention configuration reported

in the ablation study (Table~8) of Hernandez-Gutierrez et al.~\cite{Hernandez2025}

on all three datasets. The comparison is labelled as cross-study context because

the partitions and evaluation procedures are not identical.



  \item We analyse DDIM transition count as a controllable inference parameter, documenting the observed Dice--latency trade-off from 10 to 100 reverse transitions without changing the trained architecture.



  \item We present full-image and zoomed qualitative segmentation comparisons; these selected examples do not establish statistical improvement.

\end{enumerate}



\paragraph{Overview.}

Figure~\ref{fig:flowchart_architecture} illustrates the overall framework, comparing the conventional single-stage LU-Net+RA pipeline with the proposed two-stage approach, in which a diffusion module refines the initial segmentation through iterative correction of the spatially downsampled multi-channel representation.



\begin{figure}[htbp]
  \centering
  \manuscriptimage[width=0.98\linewidth]{realdata/flowchart260801-2.png}
  \caption{Workflow of the proposed two-stage framework. Stage~1 applies vessel-enhancement preprocessing and LU-Net with reverse attention to produce the coarse vessel mask $\mathbf{M}_0$. In Stage~2, the frozen mask encoder maps $\mathbf{M}_0$ to a spatially downsampled feature grid, which is noised and refined by the conditional diffusion U-Net using deterministic DDIM transitions. A separate RGB image encoder supplies the 128-dimensional conditioning vector. The frozen mask decoder produces the refined vessel-probability map $\mathbf{M}^{*}$.}
  \label{fig:flowchart_architecture}
\end{figure}



The remainder of this paper is organised as follows. Section~\ref{sec:related} reviews prior work on vessel segmentation and diffusion-based refinement. Section~\ref{sec:methods} describes the coarse predictor and proposed refinement method. Section~\ref{sec:experiments_results} presents the experimental protocol, segmentation results, and qualitative analyses. Sections~\ref{sec:discussion} and~\ref{sec:conclusion} discuss the evidence, limitations, and conclusions.

%%%%%%%%%%%%%%%%%%%%%%%%%%%%%%%%%%%%%%%%%%





%%%%%%%%%%%%%%%%%%%%%%%%%%%%%%%%%%%%%%%%%%



\section{Related Work}

\label{sec:related}



\subsection{Retinal Vessel Segmentation}



Early retinal vessel segmentation methods relied on hand-crafted features, such as Gabor filters and matched filtering, combined with thresholding or classical classifiers~\cite{Soares2006, Mendonca2006}. While effective in controlled settings, these approaches lack robustness to variations in illumination and vessel appearance.



With the advent of deep learning, convolutional neural networks (CNNs), particularly U-Net~\cite{Ronneberger2015}, became the dominant paradigm for biomedical image segmentation, achieving strong performance on benchmark datasets such as DRIVE~\cite{Staal2004}. Subsequent work has focused on improving architectural design, including nested structures (UNet++~\cite{Zhou2018UNetpp}), attention mechanisms~\cite{Oktay2018AttUNet, Schlemper2019AG}, and recurrent or residual connections~\cite{Alom2018R2UNet}, as well as incorporating vessel-specific priors such as connectivity constraints and boundary information~\cite{Fu2016DeepVessel, Liskowski2016, Shin2019Connectivity}.



More recently, lightweight architectures such as LU-Net+RA~\cite{Hernandez2025} have aimed to balance accuracy and efficiency by leveraging depth-wise separable convolutions and reverse attention. Thin-vessel recovery and structural continuity remain relevant evaluation targets, particularly in low-contrast regions. These targets motivate testing whether a separate refinement strategy changes coarse predictions beneficially under a controlled protocol.



\subsection{Diffusion Models}



Diffusion models learn to reverse a progressive noise process~\cite{Ho2020}. DDIM permits deterministic transitions along a reduced timestep sequence~\cite{Song2021}, while latent diffusion performs these updates in an autoencoder representation~\cite{Rombach2022}.



Diffusion models have also been investigated for medical-image segmentation. Wolleb et al.~\cite{Wolleb2022} generated implicit segmentation ensembles, while SegDiff and MedSegDiff adapted conditional diffusion to general and medical segmentation tasks~\cite{SegDiff,MedSegDiff}. Pre-segmentation diffusion sampling starts the reverse process from a noised preliminary mask rather than pure Gaussian noise~\cite{PDDiffusion}, which is conceptually close to the initialisation used here. A broader survey is provided by Kazerouni et al.~\cite{Kazerouni2023}.



Several diffusion methods directly address coarse-to-fine refinement. HiDiffSeg combines coarse vessel estimates, vessel skeletons, and fundus-image guidance in a hierarchical retinal segmentation model~\cite{HiDiffSeg}. RefineCatDiff uses coarse probability maps as priors in categorical diffusion~\cite{RefineCatDiff}. Li et al.~\cite{DiffusionFMD} introduce a plug-and-play diffusion denoising module for end-to-end feature-map refinement. Our diffusion refiner applies continuous Gaussian diffusion to encoded LU-Net+RA probabilities, freezes the mask autoencoder, and conditions the denoiser on a global RGB embedding. It functions as a separate post-processing stage. The parameter audit below includes its conditioning layers. These architectural differences do not establish superiority over methods that have not been evaluated in a controlled comparison.



The diffusion refiner adds noise to the encoded coarse mask and uses a decoded Dice loss to encourage agreement with the reference mask. It therefore learns Dice-regularised denoising of coarse predictions, rather than conventional diffusion of encoded ground-truth targets. Section~\ref{sec:spatial_refiner_context} reports a spatial learned-refiner experiment, and Section~\ref{sec:component_ablations} reports controlled component ablations, including target diffusion. A matched comparison with the original diffusion configuration is still needed to isolate its benefit.



\subsection{Attention Mechanisms in Medical Imaging}



Attention mechanisms are widely used in medical imaging to enhance feature representation by focusing on relevant regions. Channel and spatial attention modules, such as SE-Net~\cite{Hu2018} and CBAM~\cite{Woo2018}, improve performance by reweighting feature responses. In contrast, reverse attention~\cite{Hernandez2025} explicitly emphasises regions that are likely to be misclassified, enabling the model to focus on difficult structures.





%%%%===========================================================================================================================================

\section{Methods}

\label{sec:methods}



\subsection{Problem Formulation}

\label{sec:problem}



Retinal vessel segmentation is formulated as a dense prediction task in which each pixel of a fundus image is classified as vessel or background. Let $\mathbf{I}\in\mathbb{R}^{3\times512\times512}$ denote an RGB fundus image, and let $\mathbf{M}_{gt}\in\{0,1\}^{1\times512\times512}$ denote the corresponding ground-truth mask. The objective is to learn a mapping $\mathcal{F}:\mathbf{I}\mapsto \mathbf{M}^*$, where $\mathbf{M}^*\in[0,1]^{1\times512\times512}$ is the predicted vessel probability map (thresholded to obtain a binary mask during evaluation).



Instead of a single-stage prediction, we adopt a coarse-to-fine framework. A lightweight segmentation model first produces a coarse estimate $\mathbf{M}_0\in[0,1]^{1\times512\times512}$, which captures the main vessel structure but may contain discontinuities, missing thin vessels, and blurred boundaries.



A mask encoder maps the coarse estimate to $\mathbf{z}_0=\mathrm{Encode}(\mathbf{M}_0)$, a spatially downsampled multi-channel representation with $\mathbf{z}_0\in\mathbb{R}^{C\times H\times W}$. The conditional diffusion refiner $\mathrm{Refine}_{\theta}(\cdot)$ iteratively updates this grid using the input image $\mathbf{I}$ as conditioning information. A mask decoder $\mathrm{Decode}(\cdot)$ then maps the refined grid back to pixel space:

\begin{equation}
\mathbf{M}^* = \mathrm{Decode}\big(\mathrm{Refine}_{\theta}(\mathbf{z}_0, \mathbf{I})\big).
\end{equation}



This formulation separates coarse prediction from iterative refinement and provides a tunable additional-computation stage.



\subsection{Baseline: Lightweight U-Net with Reverse Attention}



For coarse segmentation, we use LU-Net+RA~\cite{Hernandez2025}, a lightweight encoder--decoder. The released one-channel LNet\_RA implementation used for the parameter audit employs standard convolutions and reduced channel widths. Its key component is reverse attention (RA), which emphasises uncertain regions. Given a feature map $\mathbf{F}$ and attention logits $\mathbf{A}$, the RA output is:

\begin{equation}
\mathbf{F}^{\text{RA}} = \mathbf{F}\otimes \left(1-\sigma(\mathbf{A})\right),
\end{equation}

where $\sigma(\cdot)$ is the sigmoid function and $\otimes$ denotes element-wise multiplication. Here, $\sigma(\mathbf{A})$ highlights predicted vessel regions, while its complement focuses the model on hard or misclassified areas. The resulting reverse-attention-modulated features are decoded to produce the coarse prediction $\mathbf{M}_0$ used by the diffusion refinement stage.



Reverse attention operates within the LU-Net+RA backbone and influences the coarse prediction $\mathbf{M}_0$. The diffusion refiner receives the encoded coarse mask $\mathbf{z}_0=\mathrm{Encode}(\mathbf{M}_0)$ and image embedding $\mathbf{c}=\mathrm{ImageEncoder}(\mathbf{I}_{RGB})$. It receives no separate RA map or additional RA conditioning channel.



\subsection{Proposed Conditional Diffusion Refiner}



To refine the coarse segmentation, we operate on a learned multi-channel feature grid with lower spatial resolution than the input mask. The proposed refinement framework consists of a mask autoencoder, an image encoder for conditioning, and a conditional diffusion model defined in this learned representation space.

\begin{table}[htbp]
\centering
\caption{Component interfaces and total parameter counts of the archived refinement implementation. Counts include frozen components and the separate conditioning MLP.}
\label{tab:architecture_params}
\small
\begin{tabularx}{\textwidth}{@{}p{0.24\textwidth}Xr@{}}
\toprule
Component & Interface and role & Parameters \\
\midrule
Mask autoencoder & $1\times512\times512$ mask $\rightarrow 64\times128\times128$ feature grid $\rightarrow 1\times512\times512$ probability map & 121,409 \\
Image encoder & $3\times512\times512$ RGB image $\rightarrow$ 128-dimensional global embedding & 27,712 \\
Diffusion U-Net with internal FiLM projections & Noisy feature grid and conditioning vector $\rightarrow$ predicted noise in the encoded multi-channel grid & 353,280 \\
Conditioning MLP & Two $128\rightarrow128$ linear layers applied to the image condition & 33,024 \\
\midrule
Full archived refinement module & Coarse probability map and RGB image $\rightarrow$ refined vessel-probability map & \textbf{535,425} \\
\bottomrule
\end{tabularx}
\end{table}

We count unique registered parameter tensors using \texttt{sum(p.numel() for p in model.parameters())}. The convention includes all weights and biases, including frozen pretrained parameters, while excluding non-parameter buffers, optimizer states, and duplicate aliases in checkpoint dictionaries. The released one-channel LNet\_RA baseline contains 1,946,961 parameters under this convention. Its 3,028 buffer elements are excluded. Freezing the backbone does not remove its parameters from the complete-system count.



The previously reported 502,401 refinement parameters covered the mask autoencoder (121,409), image encoder (27,712), and diffusion U-Net with internal FiLM projections (353,280). This subtotal omitted the archived implementation's separate 33,024-parameter conditioning MLP. Including it gives 535,425 refinement parameters and 2,482,386 parameters for the complete archived system, or 27.50\% above the baseline. The executable reproduction uses a different time-conditioning MLP with 24,832 parameters. Its refiner contains 527,233 parameters, and its complete system contains 2,474,194, or 27.08\% above the baseline. We report the two implementations separately because their conditioning layers differ.



\begin{table}[htbp]
\centering
\caption{Directly comparable parameter counts using the same convention. Total counts include frozen modules; the last column counts parameters updated during refinement training, when the backbone and mask autoencoder are frozen. The baseline alone has 1,946,961 trainable parameters during its own training.}
\label{tab:complete_parameter_counts}
\small
\resizebox{\textwidth}{!}{%
\begin{tabular}{lrrrr}
\toprule
Configuration & Baseline parameters & Added refiner parameters & Complete system & Updated during refinement \\
\midrule
LU-Net+RA alone & 1,946,961 & 0 & 1,946,961 & 0 \\
Archived implementation & 1,946,961 & 535,425 & 2,482,386 & 414,016 \\
New reproduction & 1,946,961 & 527,233 & 2,474,194 & 405,824 \\
\bottomrule
\end{tabular}
}
\end{table}



These counts describe parameter storage and training budgets. They do not measure computational efficiency: sequential DDIM sampling adds inference work without changing the parameter count. We checked the archived architecture against the available saved model state, but could not establish its link to the deleted original score-producing run. The executable reproduction is linked to retained checkpoints, and Table~\ref{tab:new_denoiser_layers} specifies its denoiser layer by layer.



\subsubsection{Mask Autoencoder}



The convolutional mask autoencoder maps $\mathbf{M}\in\mathbb{R}^{1\times512\times512}$ to $\mathbf{z}\in\mathbb{R}^{64\times128\times128}$, reducing each spatial dimension by $4\times$. The input contains $1\times512\times512=262{,}144$ scalar values; the encoded grid contains $64\times128\times128=1{,}048{,}576$. Thus, spatial downsampling increases the total scalar count. Parameter efficiency refers to the model budget, rather than the size of this representation. The architecture is:

\begin{itemize}

  \item Encoder: Conv($1\rightarrow32$, $3\times3$, stride 2, padding 1) + ReLU $\rightarrow$ Conv($32\rightarrow64$, $3\times3$, stride 2, padding 1) + ReLU $\rightarrow$ Conv($64\rightarrow64$, $1\times1$, stride 1).

  \item Decoder: two $4\times4$ transposed convolutions ($64\rightarrow64$ and $64\rightarrow32$; stride 2, padding 1), each followed by ReLU, then a $1\times1$ convolution ($32\rightarrow1$) and sigmoid activation.

\end{itemize}



Therefore, the spatial path is $512\times512\rightarrow256\times256\rightarrow128\times128$ in the encoder and $128\times128\rightarrow256\times256\rightarrow512\times512$ in the decoder, while channel progression is $1\rightarrow32\rightarrow64\rightarrow64$ (encoder) and $64\rightarrow64\rightarrow32\rightarrow1$ (decoder). The autoencoder has approximately $0.121$M trainable parameters in total (about $22{,}976$ in the encoder and $98{,}433$ in the decoder).



For an input ground-truth mask $\mathbf{M}_{gt}$, the autoencoder reconstruction is

\begin{equation}
\tilde{\mathbf{M}}=D_m\!\left(E_m(\mathbf{M}_{gt})\right).
\end{equation}

The reconstruction term is the mean squared error over the $N=512\times512$ pixels,

\begin{equation}
\mathcal{L}_{\text{MSE}}(\tilde{\mathbf{M}},\mathbf{M}_{gt})=
\frac{1}{N}\sum_{p=1}^{N}
\left(\tilde{M}_{p}-M_{gt,p}\right)^2.
\end{equation}

For any predicted probability map $\mathbf{P}$ and binary target $\mathbf{G}$, we use the soft Dice loss

\begin{equation}
\mathcal{L}_{\text{Dice}}(\mathbf{P},\mathbf{G})=
1-\frac{2\sum_{p=1}^{N}P_pG_p}
{\sum_{p=1}^{N}P_p+\sum_{p=1}^{N}G_p}.
\end{equation}

Accordingly, the autoencoder objective is

\begin{equation}
\mathcal{L}_{\text{AE}} =
\mathcal{L}_{\text{MSE}}(\tilde{\mathbf{M}},\mathbf{M}_{gt})+
\mathcal{L}_{\text{Dice}}(\tilde{\mathbf{M}},\mathbf{M}_{gt}).
\end{equation}

All sums are evaluated over the complete resized mask. The autoencoder is pretrained on ground-truth masks only, without additional mask augmentation. It is therefore trained on manually annotated binary masks rather than coarse LU-Net+RA probability maps. After pretraining, its parameters are frozen and used to encode $\mathbf{M}_0$. The distribution shift between binary training masks and soft coarse probability maps is analysed as a limitation in Section~\ref{sec:limitations}.



\subsubsection{Image Conditioning}



To incorporate global contextual information, the diffusion refinement stage uses a custom lightweight CNN image encoder (not ResNet/MobileNet) fed by the resized and normalised RGB fundus image. Specifically, the encoder receives $\mathbf{I}_{RGB}\in\mathbb{R}^{3\times512\times512}$ and maps it to a compact feature vector $\mathbf{c}\in\mathbb{R}^{128}$:

\begin{equation}
\mathbf{c}=\mathrm{ImageEncoder}(\mathbf{I}_{RGB}).
\end{equation}



Its layer stack is:

\begin{itemize}

  \item Conv($3\rightarrow32$, $3\times3$, stride 2, padding 1) + ReLU,

  \item Conv($32\rightarrow64$, $3\times3$, stride 2, padding 1) + ReLU,

  \item Adaptive average pooling to $1\times1$,

  \item Flatten,

  \item Linear($64\rightarrow128$) to produce the final embedding $\mathbf{c}$.

\end{itemize}



Thus, the feature map is downsampled as $512\times512\rightarrow256\times256\rightarrow128\times128\rightarrow1\times1$, then projected to a 128-dimensional conditioning vector. This image encoder has approximately $0.028$M trainable parameters ($27{,}712$ parameters) and is trained from scratch during diffusion-refiner training (no external pretrained weights are used).



The RGB conditioning branch uses different preprocessing from the coarse LU-Net+RA backbone. The backbone extracts the green channel and enhances contrast to produce $\mathbf{M}_0$. The diffusion image encoder instead receives the resized, normalised three-channel RGB fundus image. This branch supplies colour and illumination information while the denoiser operates on the encoded coarse mask.

\subsubsection{Diffusion on the Spatially Downsampled Multi-Channel Representation}



The refinement is formulated as a diffusion process on the spatially downsampled multi-channel representation. Starting from the encoded coarse mask $\mathbf{z}_0=\mathrm{Encode}(\mathbf{M}_0)$, noise is progressively added according to a cosine schedule~\cite{Nichol2021}:

\begin{equation}
\bar{\alpha}_t=\frac{\cos^2\!\left(\frac{t/T+0.008}{1.008}\cdot\frac{\pi}{2}\right)}{\cos^2\!\left(\frac{0.008}{1.008}\cdot\frac{\pi}{2}\right)},
\end{equation}

where $t\in\{0,\dots,T\}$ denotes the diffusion step and $\bar{\alpha}_t$ is the normalised cumulative signal-retention schedule, with $\bar{\alpha}_0=1$. For $t\in\{1,\dots,T\}$, the corresponding per-step quantities are $\alpha_t=\bar{\alpha}_t/\bar{\alpha}_{t-1}$ and $\beta_t=1-\alpha_t$. We use a diffusion horizon of $T=1000$ and sample training timesteps uniformly from $\{1,\dots,T-1\}$. During inference, DDIM sampling retains a shorter timestep sequence, typically containing 20 or 50 reverse transitions, as analysed in Table~\ref{tab:ddim_steps}. The forward process generates noisy feature variables:

\begin{equation}
\mathbf{z}_t=\sqrt{\bar{\alpha}_t}\mathbf{z}_0+\sqrt{1-\bar{\alpha}_t}\,\boldsymbol{\epsilon}, 
\quad \boldsymbol{\epsilon}\sim\mathcal{N}(0,\mathbf{I}),
\end{equation}

where $\boldsymbol{\epsilon}$ is Gaussian noise.



A conditional diffusion U-Net is then trained to reverse this process by predicting the noise component $\hat{\boldsymbol{\epsilon}}_\theta(\mathbf{z}_t,t,\mathbf{c})$, conditioned on both the noisy feature tensor $\mathbf{z}_t$ and the image embedding $\mathbf{c}$.





\subsubsection{Feature-Conditioning Injection}



Learned conditioning layers project the 128-dimensional image embedding $\mathbf{c}$ into the diffusion U-Net's intermediate residual features. The timestep embedding enters through the same residual-block mechanism. This conditioning supplies global image context throughout denoising without a spatial cross-attention map. Reverse attention is confined to the LU-Net+RA backbone that produces $\mathbf{M}_0$; the denoiser receives no separate RA tensor.



\subsubsection{Training Objective}



The model is trained to predict the noise added during the forward diffusion process while incorporating segmentation-aware constraints. Let $N_z$ be the number of elements in a feature tensor. The mean DDPM noise-prediction loss is

\begin{equation}
\mathcal{L}_{\text{DDPM}}=
\mathbb{E}_{\mathbf{z}_0,t,\boldsymbol{\epsilon}}\left[
\frac{1}{N_z}\left\|
\hat{\boldsymbol{\epsilon}}_\theta(\mathbf{z}_t,t,\mathbf{c})-
\boldsymbol{\epsilon}
\right\|_2^2
\right],
\end{equation}

with $\boldsymbol{\epsilon}\sim\mathcal{N}(0,\mathbf{I})$ denoting Gaussian noise and $\mathbf{z}_t$ the noisy feature tensor obtained from $\mathbf{z}_0$ at timestep $t$.

The noise prediction is then used to estimate the clean feature tensor corresponding to the sampled timestep:

\begin{equation}
\hat{\mathbf{z}}_0^{(t)}=
\frac{\mathbf{z}_t-
\sqrt{1-\bar{\alpha}_t}\,
\hat{\boldsymbol{\epsilon}}_\theta(\mathbf{z}_t,t,\mathbf{c})}
{\sqrt{\bar{\alpha}_t}}.
\end{equation}

Because the pretrained mask decoder is frozen and includes its output sigmoid, the associated vessel-probability map is

\begin{equation}
\hat{\mathbf{M}}_t=D_m\!\left(\hat{\mathbf{z}}_0^{(t)}\right).
\end{equation}

The segmentation-aware term is therefore $\mathcal{L}_{\text{Dice}}(\hat{\mathbf{M}}_t,\mathbf{M}_{gt})$, using the pixel-space definition given above, and the complete diffusion-refiner objective is

\begin{equation}
\mathcal{L}=
\mathcal{L}_{\text{DDPM}}+
\lambda_{\text{dice}}
\mathcal{L}_{\text{Dice}}(\hat{\mathbf{M}}_t,\mathbf{M}_{gt}).
\end{equation}

In the implementation evaluated here, $\lambda_{\text{dice}}=0.5$; no separate Sobel edge-loss term is applied. A supplementary comparison with weights 0, 0.25, and 1 is reported in Section~\ref{sec:dice_weight_ablation}.



The objective combines two distinct learning signals. The DDPM term predicts Gaussian noise added to $E_m(\mathbf{M}_0)$, encouraging recovery of the encoded coarse mask. The decoded Dice term encourages agreement with the ground-truth vessel mask. This is Dice-regularised coarse-mask denoising, rather than unconditional generation of $E_m(\mathbf{M}_{gt})$ or target diffusion with a separate coarse-mask condition.



The models are trained in two stages. The mask autoencoder is pretrained for 50 epochs using AdamW with a learning rate of $1\times10^{-4}$ on ground-truth vessel masks only. Its weights are then frozen, and the diffusion refiner is trained for 100 epochs using AdamW at $1\times10^{-4}$ with cosine annealing over the full training schedule. During diffusion training, only the image encoder, conditioning multilayer perceptron (MLP), and diffusion U-Net are updated. A classifier-free-guidance conditioning dropout of 10\% is applied during training. No additional image or mask augmentation is used. Images and masks are resized to $512\times512$, and LU-Net+RA probability maps are loaded as the coarse conditions. No explicit field-of-view loss mask or crop is applied during refiner training.



\subsubsection{DDIM Sampling}



For inference, we employ DDIM sampling~\cite{Song2021} to generate a refined feature representation from the initial code $\mathbf{z}_0=\mathrm{Encode}(\mathbf{M}_0)$. Instead of starting from pure noise, the sampling process is initialised from a noisy version of $\mathbf{z}_0$, preserving an explicit connection to the coarse prediction throughout the refinement trajectory.



\begin{table}[htbp]
\caption{Sampling procedure for refinement of the encoded multi-channel representation.}
\label{alg:latent_refinement}
\centering
\begin{minipage}{0.96\textwidth}
\small
\textbf{Inputs:} coarse mask $\mathbf{M}_0$; RGB image $\mathbf{I}_{RGB}$; encoders $E_m,E_I$; decoder $D_m$; noise network $\boldsymbol{\epsilon}_\theta$; schedule $\{\bar{\alpha}_t\}_{t=0}^{T}$; refinement strength $\rho$; DDIM transitions $S$; guidance scale $s$; sample count $K$.\par
\begin{enumerate}[leftmargin=2em,itemsep=2pt,topsep=4pt]
  \item Encode $\mathbf{z}_0=E_m(\mathbf{M}_0)$ and $\mathbf{c}=E_I(\mathbf{I}_{RGB})$; set the null condition $\mathbf{c}_{\emptyset}=\mathbf{0}$.
  \item Set $t_s=\lfloor\rho T\rfloor$ and $t_i=\lfloor t_s(S-i)/S\rfloor$ for $i=0,\ldots,S$, so $t_0=t_s$ and $t_S=0$.
  \item For each sample $k=1,\ldots,K$, draw $\boldsymbol{\epsilon}^{(k)}\sim\mathcal{N}(0,\mathbf{I})$ and initialise $\mathbf{z}^{(k)}_{t_0}=\sqrt{\bar{\alpha}_{t_0}}\mathbf{z}_0+\sqrt{1-\bar{\alpha}_{t_0}}\boldsymbol{\epsilon}^{(k)}$.
  \item For $i=0,\ldots,S-1$, evaluate $\boldsymbol{\epsilon}_{c}=\boldsymbol{\epsilon}_\theta(\mathbf{z}^{(k)}_{t_i},t_i,\mathbf{c})$ and $\boldsymbol{\epsilon}_{u}=\boldsymbol{\epsilon}_\theta(\mathbf{z}^{(k)}_{t_i},t_i,\mathbf{c}_{\emptyset})$, then form $\hat{\boldsymbol{\epsilon}}=\boldsymbol{\epsilon}_{u}+s(\boldsymbol{\epsilon}_{c}-\boldsymbol{\epsilon}_{u})$.
  \item Estimate $\hat{\mathbf{z}}_0=(\mathbf{z}^{(k)}_{t_i}-\sqrt{1-\bar{\alpha}_{t_i}}\hat{\boldsymbol{\epsilon}})/\sqrt{\bar{\alpha}_{t_i}}$ and update $\mathbf{z}^{(k)}_{t_{i+1}}=\sqrt{\bar{\alpha}_{t_{i+1}}}\hat{\mathbf{z}}_0+\sqrt{1-\bar{\alpha}_{t_{i+1}}}\hat{\boldsymbol{\epsilon}}$.
  \item Decode $\tilde{\mathbf{M}}^{(k)}=D_m(\mathbf{z}^{(k)}_{t_S})$ and return $\mathbf{M}^*=K^{-1}\sum_{k=1}^{K}\tilde{\mathbf{M}}^{(k)}$.
\end{enumerate}
\end{minipage}
\end{table}



The starting timestep $t_s$ sets the refinement strength. Larger values add more noise and permit greater correction; smaller values preserve more of the LU-Net+RA prediction. Uniform spacing of $S+1$ timesteps between $t_s$ and zero gives exactly $S$ reverse transitions. Each transition evaluates the diffusion U-Net with the image condition and with a null condition, then combines the predictions through classifier-free guidance. A deterministic DDIM update ($\eta_{\mathrm{DDIM}}=0$) advances the feature tensor to the next timestep. The decoder already ends with a sigmoid, so Table~\ref{alg:latent_refinement} applies no second sigmoid. The primary segmentation results use $K=1$. A reverse trajectory is deterministic for a fixed initial noise sample, but the complete procedure varies with that sample.



%%%%===========================================================================================================================================

%%%%===========================================================================================================================================



\section{Experiments and Results}

\label{sec:experiments_results}



\subsection{Datasets and Preprocessing}



We evaluate the refiner on three public retinal vessel datasets: the Child Heart and Health Study in England Database 1 (CHASE-DB1)~\cite{CHASEDB1}, Digital Retinal Images for Vessel Extraction (DRIVE)~\cite{Staal2004}, and High-Resolution Fundus (HRF)~\cite{Odstrcilik2013HRF}. Table~\ref{tab:datasets} reports the released image sizes used by the preprocessing pipeline and the study partitions stated for the experiments.



\begin{table}[htbp]
\centering
\caption{Dataset composition and partitions stated for this study. Image sizes refer to the released or cropped images used by the pipeline, not necessarily the acquisition resolution.}
\label{tab:datasets}
\small
\begin{tabularx}{\textwidth}{@{}l c c X c@{}}
\toprule
Dataset & Images & Released Image Size & Study Partition & Test $n$ \\
\midrule
CHASE-DB1 & 28 & $999\times960$ & 20 train / 8 test & 8 \\
DRIVE & 40 & $565\times584$ & 20 train / 20 test & 20 \\
HRF & 45 & $3504\times2336$ & 33 train / 12 test & 12 \\
\bottomrule
\end{tabularx}
\end{table}



CHASE-DB1 contains images of both eyes from 14 children. The first-observer annotation is used in the qualitative figures; the observer used for all quantitative calculations and subject-level independence of the stated 20/8 partition are not documented in the available experiment files. DRIVE has an official 20-image training set and 20-image test set. HRF contains 15 healthy, 15 diabetic-retinopathy, and 15 glaucomatous images; whether the stated 33/12 partition was stratified by pathology is not documented. These unresolved protocol details limit exact replication but do not change the dataset counts reported in Table~\ref{tab:datasets}.



Two preprocessing streams are used. The LU-Net+RA backbone follows the preprocessing pipeline of Hernandez-Gutierrez et al.~\cite{Hernandez2025}: green-channel extraction, contrast-limited adaptive histogram equalisation (CLAHE), inverse gamma correction, resizing to $512\times512$, and normalisation to $[0,1]$. This stream produces the coarse probability map $\mathbf{M}_0$. The refiner receives $\mathbf{M}_0$, while its image encoder receives the corresponding resized and normalised RGB fundus image. The 20- and 50-transition variants use the same stated partitions and coarse inputs.



\subsection{Training and Implementation Details}



Training is sequential. LU-Net+RA is trained first and then frozen. Its stored probability maps condition the refiner, whose training does not update the backbone. The coarse maps used to train the refiner were generated from images that had also been used to fit LU-Net+RA; they are therefore in-sample rather than out-of-fold predictions. Freezing prevents joint adaptation between the stages, but it does not remove the possible train-test shift between in-sample coarse predictions and predictions for unseen images.



The mask autoencoder is pretrained for 50 epochs on the corresponding training masks and then frozen. The image encoder, conditioning layers, and diffusion U-Net are trained for 100 epochs using AdamW, an initial learning rate of $10^{-4}$, and cosine annealing. Conditioning dropout is 10\%, and no additional geometric, image, or mask augmentation is applied. Images, masks, and coarse probability maps are resized deterministically to $512\times512$. The original files do not specify interpolation kernels, aspect-ratio handling, training batch size, AdamW weight decay, checkpoint selection, or a validation partition.



Refinement-stage timing uses a batch size of 1 on one NVIDIA GeForce RTX 4070 Ti SUPER graphics processing unit, evaluation mode, and disabled gradient computation. The timing includes mask encoding, DDIM refinement, and mask decoding but excludes the LU-Net+RA forward pass. Mixed-precision inference was not used. Because the warm-up count, number of repetitions, and timing dispersion are unavailable, all runtime values are reported as approximate rather than as benchmark estimates.



\subsection{Evaluation Protocol and Metrics}

\label{sec:metrics}



Vessel pixels are treated as the positive class. Let TP, TN, FP, and FN denote true-positive, true-negative, false-positive, and false-negative pixel counts, respectively. The metrics appearing in the result tables are

\begin{align}
\mathrm{Dice} &= \frac{2\mathrm{TP}}{2\mathrm{TP}+\mathrm{FP}+\mathrm{FN}}, \\
\mathrm{Accuracy} &= \frac{\mathrm{TP}+\mathrm{TN}}
{\mathrm{TP}+\mathrm{TN}+\mathrm{FP}+\mathrm{FN}}, \\
\mathrm{Sensitivity} &= \frac{\mathrm{TP}}{\mathrm{TP}+\mathrm{FN}}, \\
\mathrm{Specificity} &= \frac{\mathrm{TN}}{\mathrm{TN}+\mathrm{FP}}.
\end{align}

The reported metrics were restricted to a field-of-view (FOV) mask. Training losses, by contrast, were evaluated over the full resized canvas, with pixels outside the retinal field represented as background. The authors confirmed that the baseline and refinement results for each dataset used the same test images, coarse-predictor checkpoint, reference annotations, FOV masks, probability threshold, and metric aggregation rule. This equivalence is author-confirmed and has not been independently checked against the original per-image diffusion outputs and checkpoint records. CHASE-DB1 Dice is reported as the mean across eight test images; the documented DRIVE and HRF baseline evaluations used a threshold of 0.5 and per-image averaging. The baseline details are given below. These results are single-run point estimates and do not establish statistical significance.



\subsection{Quantitative Segmentation Results}



Table~\ref{tab:segmentation_results} reports the LU-Net+RA coarse baselines and the diffusion-refinement configurations on the stated partitions. Under the author-confirmed matching of test images, coarse-predictor checkpoints, and evaluation settings, the 50-transition configuration increased mean Dice from 0.8001 to 0.835 on CHASE-DB1, from 0.7879 to 0.835 on DRIVE, and from 0.7325 to 0.793 on HRF. The absolute increases were 0.0349, 0.0471, and 0.0605 (3.49, 4.71, and 6.05 percentage points), respectively. The 20-transition configuration also exceeded the corresponding within-study Dice baseline on all three partitions. These mean scores do not establish statistical significance. On DRIVE, sensitivity at 50 transitions was 0.825, below the baseline value of 0.8602, despite higher Dice; sensitivity also decreased from 0.835 at 20 transitions to 0.825 at 50 transitions. Thus, higher Dice does not imply improvement in every metric.



\begin{table}[htbp]
\centering
\caption{Test-set results from the rerun LU-Net+RA baseline and the diffusion-refinement configurations on eight CHASE-DB1, 20 DRIVE, and 12 HRF images. Matching of test images, coarse-predictor checkpoints, annotations, and evaluation settings is author-confirmed, without independent verification against the original diffusion outputs. Values are single-run point estimates.}
\label{tab:segmentation_results}
\small
\resizebox{\textwidth}{!}{%
\begin{tabular}{llccccc}
\toprule
Dataset & Configuration & Test $n$ & Dice & Accuracy & Sensitivity & Specificity \\
\midrule
CHASE-DB1
& LU-Net+RA (same protocol) & 8 & 0.8001 & 0.932 & 0.801 & .0911 \\
& Diffusion refinement (20) & 8 & 0.825 & 0.967 & 0.858 & 0.982 \\
& Diffusion refinement (50) & 8 & 0.835 & 0.968 & 0.868 & 0.984 \\
\midrule
DRIVE
& LU-Net+RA (same protocol) & 20 & 0.7879 & 0.9413 & 0.8602 & 0.9535 \\
& Diffusion refinement (20) & 20 & 0.795 & 0.965 & 0.835 & 0.980 \\
& Diffusion refinement (50) & 20 & 0.835 & 0.967 & 0.825 & 0.982 \\
\midrule
HRF
& LU-Net+RA (same protocol) & 12 & 0.7325 & 0.9480 & 0.7944 & 0.9631 \\
& Diffusion refinement (20) & 12 & 0.762 & 0.963 & 0.788 & 0.977 \\
& Diffusion refinement (50) & 12 & 0.793 & 0.966 & 0.815 & 0.980 \\
\bottomrule
\end{tabular}
}
\end{table}





\subsection{DRIVE False-Positive and False-Negative Rate Trade-off}



The reported DRIVE metrics do not change monotonically with the DDIM transition count. From 20 to 50 transitions, Dice increases from 0.795 to 0.835 (+0.040), while sensitivity decreases from 0.835 to 0.825 (-0.010) and specificity increases from 0.980 to 0.982 (+0.002). To express this trade-off quantitatively, the false-negative and false-positive rates are

\begin{equation}
\mathrm{FNR}=\frac{\mathrm{FN}}{\mathrm{TP}+\mathrm{FN}}=1-\mathrm{Sensitivity},\qquad
\mathrm{FPR}=\frac{\mathrm{FP}}{\mathrm{TN}+\mathrm{FP}}=1-\mathrm{Specificity}.
\end{equation}

Thus, the reported values imply an FNR increase from 0.165 to 0.175 (1.0 percentage point, approximately 6.1\% relative) and an FPR decrease from 0.020 to 0.018 (0.2 percentage points, approximately 10\% relative). These rates are derived from rounded summary metrics; they are not newly measured pixel counts. Under identical test images, annotations, FOV masks, thresholds, and aggregation, their direction is consistent with a shift toward fewer false-positive predictions at the cost of reduced vessel recall. The rates have different denominators and cannot be directly compared as absolute numbers of pixels.



The higher Dice therefore does not establish uniformly better vessel detection. Nor do these aggregate rates establish that false positives decreased or false negatives increased in every image, or that the missed pixels belong specifically to thin vessels. Original paired prediction maps and per-image TP, TN, FP, and FN counts are unavailable for independent verification of this comparison. Paired confusion-count analysis, threshold sweeps, and vessel-width-stratified evaluation are needed to quantify absolute error changes and determine which vessel structures are affected.



\subsection{Baseline Evaluation Details and Dice Differences}



The LU-Net+RA baseline was rerun and evaluated on the test partition of each dataset. The resulting test-set mean Dice scores were 0.8001 on the eight CHASE-DB1 test images, 0.7879 on the 20 DRIVE test images, and 0.7325 on the 12 HRF test images. These are measured test results from the within-study baseline evaluations, not values taken from the reference publication. The CHASE-DB1 result was supplied by the authors from their test run; the DRIVE and HRF results were recorded in the new baseline evaluations described below.



The DRIVE baseline evaluation used the authors' released LNet\_RA architecture with one green input channel. The 20 training images were divided into 16 training and four validation images using seed 42; the 20 official test images were reserved for evaluation. The local annotations were stored as \texttt{manual1} PNG files, whose original acquisition provenance remains undocumented. Evaluation was restricted to the supplied DRIVE FOV masks, resized with nearest-neighbour interpolation. Images were resized to $512\times512$, the probability threshold was fixed at 0.5, and metrics were averaged per image. The model used AdamW with learning rate and weight decay of 0.001, Dice loss, batch size four, and validation-based early stopping with patience 20. One training seed was evaluated.



The HRF baseline evaluation used 27 training, six validation, and 12 test images, stratified across healthy, diabetic-retinopathy, and glaucomatous groups with seed 42. Vessel annotations and FOV masks were decoded from the local annotation files. Evaluation used resized $512\times512$ inputs, a fixed probability threshold of 0.5, and per-image averaging within the FOV. These baseline records document the new evaluations; matching them to the earlier refinement results rests on the authors' confirmation rather than an independent audit of the original diffusion outputs.



\begin{table}[htbp]
\centering
\caption{Within-study Dice differences based on author-confirmed matching of baseline and refinement evaluations. Differences are absolute point-estimate changes, not estimates of statistical significance.}
\label{tab:baseline_comparison}
\begin{tabular}{lccccc}
\toprule
Dataset & Baseline & 20 transitions & Difference & 50 transitions & Difference \\
\midrule
CHASE-DB1 & 0.8001 & 0.825 & +0.0249 & 0.835 & +0.0349 \\
DRIVE & 0.7879 & 0.795 & +0.0071 & 0.835 & +0.0471 \\
HRF & 0.7325 & 0.762 & +0.0295 & 0.793 & +0.0605 \\
\bottomrule
\end{tabular}
\end{table}



The comparisons in Table~\ref{tab:baseline_comparison} use within-study coarse baseline scores rather than published scores from different protocols. They favour refinement in mean Dice on the stated partitions, conditional on the author-confirmed matching. Per-image paired differences, confidence intervals, and repeated training runs are still needed to assess robustness. The single training seed and undocumented acquisition provenance of the local DRIVE annotations remain limitations.



\subsection{Published Literature Context}



Table~\ref{tab:published_context} separately reports the baseline U-Net and selected lightweight U-Net with reverse attention results from Table~8 of Hernandez-Gutierrez et al.~\cite{Hernandez2025}. The selected LU-Net+RA row uses AdamW, Dice loss without binary cross-entropy loss, and reverse attention on each dataset. This single, explicit source avoids mixing values from the abstract, dataset-comparison tables, discussion, and conclusion of that article, which report internally different CHASE-DB1 and HRF values.



\begin{table}[htbp]
\centering
\caption{Published contextual results from Table~8 of Hernandez-Gutierrez et al.~\cite{Hernandez2025}. These values were obtained under that study's partitioning, validation, preprocessing, and metric procedures and are not controlled before-refinement measurements for Table~\ref{tab:segmentation_results}.}
\label{tab:published_context}
\small
\resizebox{\textwidth}{!}{%
\begin{tabular}{llcccc}
\toprule
Dataset & Published Configuration & Dice & Accuracy & Sensitivity & Specificity \\
\midrule
CHASE-DB1
& Baseline U-Net & 0.7830 & 0.9725 & 0.7702 & 0.9864 \\
& LU-Net+RA (selected Table~8 row) & 0.7946 & 0.9739 & 0.7795 & 0.9873 \\
\midrule
DRIVE
& Baseline U-Net & 0.7373 & 0.8931 & 0.8687 & 0.9515 \\
& LU-Net+RA (selected Table~8 row) & 0.7871 & 0.9113 & 0.7421 & 0.9837 \\
\midrule
HRF
& Baseline U-Net & 0.6417 & 0.8710 & 0.8559 & 0.9531 \\
& LU-Net+RA (selected Table~8 row) & 0.7756 & 0.8474 & 0.7794 & 0.9793 \\
\bottomrule
\end{tabular}
}
\end{table}



The 50-transition Dice estimates exceed the selected published LU-Net+RA values. The 20-transition HRF estimate (0.762), however, is below the published value (0.7756). These studies use different cases and evaluation pipelines, so their differences provide literature context rather than refinement effect estimates. The HRF score at 20 transitions exceeds the within-study baseline of 0.7325 by 0.0295. Within-study comparisons use Table~\ref{tab:segmentation_results} and the author-confirmed protocol matching described above.







\subsection{Paired Results, Repeated Seeds, and Small-Test-Set Uncertainty}

\label{sec:per_image_uncertainty}



We analyse the current author-supplied evaluation table containing 5000 records: 1000 for CHASE-DB1, 2500 for DRIVE, and 1500 for HRF. Each test image has five labels in the training-seed column, five sampling-seed labels, and five transition settings (10, 20, 30, 50, and 100). Both seed columns use 42, 123, 2026, 3000, and 3001. The table explicitly separates these two levels; sampling realisations are not treated as independent training runs or independent test images. The statistical analysis treats the training-seed labels as paired run units, conditional on their representing independent training executions; the associated execution records were not independently inspected.



The authors confirm that the coarse and refined predictions used the same test images, ground-truth masks, preprocessing, image resolution, field-of-view masks, binarisation thresholds, and Dice calculation. The refiner received the same coarse predictions used to calculate the paired baseline metrics. All five segmentation metrics and paired differences in the current table are numerically consistent with its TP, FP, FN, and TN counts. This is a numerical check of the supplied tabulation, rather than independent verification against prediction masks.



This analysis supersedes the earlier supplementary paired tabulations in this section and must be distinguished from the original single-run dataset and transition-count tables elsewhere in the manuscript. It also differs from the separately executed local reproduction. In particular, its baseline means are 0.810514 on CHASE-DB1, 0.791814 on DRIVE, and 0.732121 on HRF; these should not be substituted into comparisons with unrelated predictions.



We first average the five sampling realisations within each image and training label. Each training-label mean then averages the fixed test images. Table~\ref{tab:paired_five_seed_means} reports refined Dice at 50 transitions by training label, and Table~\ref{tab:uploaded_step_mean_sd} reports means and sample standard deviations across the five training-label means at every transition setting. These SDs describe the supplied run-level variation on a fixed test set, not uncertainty across patients.



\begin{table}[htbp]
\centering
\caption{Refined mean Dice at 50 transitions by training-seed label, after averaging sampling realisations and test images.}
\label{tab:paired_five_seed_means}
\scriptsize
\begin{tabular}{lrrrrr}
\toprule
Dataset & 42 & 123 & 2026 & 3000 & 3001 \\
\midrule
CHASE-DB1 & 0.834626923 & 0.832644216 & 0.834060111 & 0.833626098 & 0.840040756 \\
DRIVE & 0.842266236 & 0.845773452 & 0.832293158 & 0.829758676 & 0.824909831 \\
HRF & 0.786260082 & 0.785510915 & 0.796018320 & 0.794523938 & 0.802684453 \\
\bottomrule
\end{tabular}
\end{table}

\begin{table}[htbp]
\centering
\caption{Current supplied-table mean Dice and SD across five training-label means. The five sampling realisations are averaged within each run. Calculations retain full precision.}
\label{tab:uploaded_step_mean_sd}
\scriptsize
\begin{tabular}{lrrrrr}
\toprule
Dataset & Steps & Baseline mean & Baseline SD & Refined mean & Refined SD \\
\midrule
CHASE-DB1 & 10 & 0.810514023 & 0.001805530 & 0.834993547 & 0.002891775 \\
CHASE-DB1 & 20 & 0.810514023 & 0.001805530 & 0.835000734 & 0.002907682 \\
CHASE-DB1 & 30 & 0.810514023 & 0.001805530 & 0.835004217 & 0.002907204 \\
CHASE-DB1 & 50 & 0.810514023 & 0.001805530 & 0.834999621 & 0.002909859 \\
CHASE-DB1 & 100 & 0.810514023 & 0.001805530 & 0.835000101 & 0.002907811 \\
DRIVE & 10 & 0.791814385 & 0.007557284 & 0.835003891 & 0.008735743 \\
DRIVE & 20 & 0.791814385 & 0.007557284 & 0.834999788 & 0.008739377 \\
DRIVE & 30 & 0.791814385 & 0.007557284 & 0.834999933 & 0.008738995 \\
DRIVE & 50 & 0.791814385 & 0.007557284 & 0.835000270 & 0.008738928 \\
DRIVE & 100 & 0.791814385 & 0.007557284 & 0.835000324 & 0.008738913 \\
HRF & 10 & 0.732121014 & 0.003489917 & 0.792987471 & 0.007220347 \\
HRF & 20 & 0.732121014 & 0.003489917 & 0.792997823 & 0.007195086 \\
HRF & 30 & 0.732121014 & 0.003489917 & 0.793004558 & 0.007190764 \\
HRF & 50 & 0.732121014 & 0.003489917 & 0.792999542 & 0.007189127 \\
HRF & 100 & 0.732121014 & 0.003489917 & 0.793000443 & 0.007186992 \\
\bottomrule
\end{tabular}
\end{table}



We test three planned comparisons per dataset: baseline versus 50 transitions, 10 versus 50, and 50 versus 100. Two-sided paired $t$ tests and exact Wilcoxon signed-rank tests use the five training-label means. For signed-rank tests, we enumerate all signs, exclude zero differences, and average tied absolute ranks. Holm correction is applied separately to each test type across the nine comparisons. The $t$ test assumes approximately normal independent run-level differences; the signed-rank test assumes independent units and symmetric differences. Five runs provide little evidence for assessing these assumptions, so both analyses are reported.



\begin{table}[htbp]
\centering
\caption{Paired comparisons across five training-label means. Differences are second setting minus first. The two $p$ columns are separately Holm-adjusted over nine comparisons per test type.}
\label{tab:uploaded_training_paired_tests}
\scriptsize
\begin{tabular}{llrrr}
\toprule
Dataset & Comparison & Mean difference & Holm $p_t$ & Holm $p_W$ \\
\midrule
CHASE-DB1 & Baseline to 50 & 0.024485598 & 0.001365 & 0.562500 \\
CHASE-DB1 & 10 to 50 & 0.000006074 & 1.000000 & 1.000000 \\
CHASE-DB1 & 50 to 100 & 0.000000480 & 1.000000 & 1.000000 \\
DRIVE & Baseline to 50 & 0.043185885 & $1.728\times10^{-6}$ & 0.562500 \\
DRIVE & 10 to 50 & -0.000003621 & 1.000000 & 1.000000 \\
DRIVE & 50 to 100 & 0.000000053 & 1.000000 & 1.000000 \\
HRF & Baseline to 50 & 0.060878527 & 0.001043 & 0.562500 \\
HRF & 10 to 50 & 0.000012071 & 1.000000 & 1.000000 \\
HRF & 50 to 100 & 0.000000902 & 1.000000 & 1.000000 \\
\bottomrule
\end{tabular}
\end{table}



Baseline--refinement differences are positive on all three datasets and have Holm-adjusted paired $t$ probabilities below 0.05, conditional on the stated run-level assumptions. However, all three corresponding exact signed-rank probabilities are 0.0625 before adjustment and 0.5625 after adjustment. With five nonzero untied pairs, 0.0625 is the smallest attainable two-sided exact signed-rank probability. Thus the parametric and nonparametric analyses do not support identical significance claims.



None of the 10-versus-50 or 50-versus-100 comparisons reaches 0.05 in either adjusted test. The 50-to-100 mean changes are only 0.000000480 for CHASE-DB1, 0.000000053 for DRIVE, and 0.000000902 for HRF. These current-table results do not establish a reliable benefit of additional transitions or equivalence between settings. They do not explain the original single-run increase from 0.835 to 0.836, which remains a distinct result and may reflect sampling variation.



To assess test-unit variation, we additionally average all training and sampling realisations within each image. CHASE testing averages both eyes within each of four subjects; DRIVE and HRF use twenty and twelve image units. Patient independence cannot be verified for DRIVE or HRF from the supplied identifiers. Secondary two-sided tests use separate Holm correction across the three datasets for each test type; these tests answer a different question from the run-level comparisons.



\begin{table}[htbp]
\centering
\caption{Secondary subject/image-level paired tests at 50 transitions. Both test families are separately Holm-adjusted across three datasets. CHASE units are subjects; DRIVE and HRF units are images with unverified patient independence.}
\label{tab:paired_statistical_tests}
\scriptsize
\begin{tabular}{lrrr}
\toprule
Dataset & Units & Holm $p_t$ & Holm $p_W$ \\
\midrule
CHASE-DB1 & 4 & 0.098927 & 0.125000 \\
DRIVE & 20 & $1.510\times10^{-29}$ & $5.722\times10^{-6}$ \\
HRF & 12 & $7.980\times10^{-10}$ & 0.000977 \\
\bottomrule
\end{tabular}
\end{table}



The CHASE subject-level comparison does not reach 0.05 in either test ($p_t=0.098927$ and $p_W=0.125$ after adjustment). DRIVE and HRF reach this threshold under the stated image-level assumptions. These conditional findings should not be described as broad patient-level robustness.



Exploratory 95\% percentile-bootstrap intervals use 50,000 draws with seed 20261008. We average sampling realisations, then jointly resample training labels and test units. CHASE resampling retains both eyes of each subject together. DRIVE uses image resampling, and HRF uses image resampling within pathology strata. Table~\ref{tab:per_image_paired_intervals} reports the baseline--refinement differences at 50 transitions. With five training labels and four CHASE subjects, interval coverage may be unreliable. Positive bootstrap intervals do not override the nonsignificant subject-level tests.



\begin{table}[htbp]
\centering
\caption{Current supplied-table paired mean Dice at 50 transitions and exploratory crossed training-label/test-unit bootstrap intervals for the difference.}
\label{tab:per_image_paired_intervals}
\scriptsize
\begin{tabular}{lrrrl}
\toprule
Dataset & Baseline & Refined & Difference & 95\% interval \\
\midrule
CHASE-DB1 & 0.810514 & 0.835000 & +0.024486 & $[0.006382,\,0.042746]$ \\
DRIVE & 0.791814 & 0.835000 & +0.043186 & $[0.041895,\,0.044248]$ \\
HRF & 0.732121 & 0.793000 & +0.060879 & $[0.052088,\,0.069338]$ \\
\bottomrule
\end{tabular}
\end{table}

\begin{table}[htbp]
\centering
\caption{CHASE-DB1 current supplied-table per-image means at 50 transitions, averaging 25 training/sampling realisations per image. Values are rounded upward to three decimals for display; statistical calculations use unrounded values.}
\label{tab:chase_per_image_corrected}
\scriptsize
\begin{tabular}{lrrr}
\toprule
Image & Baseline Dice & Refined Dice & Difference \\
\midrule
Image\_02L & 0.819 & 0.863 & 0.045 \\
Image\_02R & 0.813 & 0.841 & 0.029 \\
Image\_05L & 0.839 & 0.893 & 0.054 \\
Image\_05R & 0.845 & 0.887 & 0.042 \\
Image\_09L & 0.806 & 0.819 & 0.013 \\
Image\_09R & 0.797 & 0.800 & 0.003 \\
Image\_13L & 0.792 & 0.807 & 0.015 \\
Image\_13R & 0.775 & 0.773 & -0.002 \\
\bottomrule
\end{tabular}
\end{table}

\begin{table}[htbp]
\centering
\caption{DRIVE current supplied-table per-image means at 50 transitions, averaging 25 training/sampling realisations per image. Values are rounded upward to three decimals for display; statistical calculations use unrounded values.}
\label{tab:drive_per_image_supplied}
\scriptsize
\begin{tabular}{lrrr}
\toprule
Image & Baseline Dice & Refined Dice & Difference \\
\midrule
01\_test & 0.798 & 0.840 & 0.043 \\
02\_test & 0.828 & 0.873 & 0.046 \\
03\_test & 0.788 & 0.832 & 0.045 \\
04\_test & 0.805 & 0.849 & 0.045 \\
05\_test & 0.797 & 0.842 & 0.045 \\
06\_test & 0.786 & 0.830 & 0.045 \\
07\_test & 0.801 & 0.845 & 0.045 \\
08\_test & 0.781 & 0.826 & 0.045 \\
09\_test & 0.784 & 0.827 & 0.044 \\
10\_test & 0.786 & 0.829 & 0.043 \\
11\_test & 0.774 & 0.815 & 0.041 \\
12\_test & 0.797 & 0.841 & 0.044 \\
13\_test & 0.794 & 0.837 & 0.044 \\
14\_test & 0.790 & 0.833 & 0.043 \\
15\_test & 0.764 & 0.805 & 0.041 \\
16\_test & 0.809 & 0.854 & 0.045 \\
17\_test & 0.790 & 0.834 & 0.044 \\
18\_test & 0.787 & 0.829 & 0.043 \\
19\_test & 0.803 & 0.843 & 0.041 \\
20\_test & 0.784 & 0.827 & 0.043 \\
\bottomrule
\end{tabular}
\end{table}

\begin{table}[htbp]
\centering
\caption{HRF current supplied-table per-image means at 50 transitions, averaging 25 training/sampling realisations per image. Values are rounded upward to three decimals for display; statistical calculations use unrounded values.}
\label{tab:hrf_per_image_supplied}
\scriptsize
\begin{tabular}{lrrr}
\toprule
Image & Baseline Dice & Refined Dice & Difference \\
\midrule
01\_g & 0.711 & 0.777 & 0.067 \\
02\_h & 0.799 & 0.872 & 0.074 \\
04\_g & 0.742 & 0.809 & 0.068 \\
07\_dr & 0.734 & 0.794 & 0.060 \\
09\_h & 0.751 & 0.831 & 0.080 \\
10\_dr & 0.730 & 0.779 & 0.050 \\
12\_dr & 0.698 & 0.758 & 0.061 \\
13\_g & 0.704 & 0.750 & 0.046 \\
13\_h & 0.776 & 0.844 & 0.068 \\
14\_g & 0.702 & 0.754 & 0.053 \\
14\_h & 0.763 & 0.818 & 0.056 \\
15\_dr & 0.681 & 0.734 & 0.054 \\
\bottomrule
\end{tabular}
\end{table}

\begin{table}[htbp]
\centering
\caption{Mean segmentation metrics at 50 transitions from the current supplied table, rounded upward to three decimals for display.}
\label{tab:supplied_paired_metric_means}
\scriptsize
\begin{tabular}{llrrrrr}
\toprule
Dataset & Method & Dice & IoU & Sensitivity & Specificity & Accuracy \\
\midrule
CHASE-DB1 & Coarse baseline & 0.811 & 0.682 & 0.807 & 0.987 & 0.975 \\
CHASE-DB1 & Refinement & 0.835 & 0.719 & 0.830 & 0.989 & 0.979 \\
DRIVE & Coarse baseline & 0.792 & 0.656 & 0.827 & 0.963 & 0.945 \\
DRIVE & Refinement & 0.836 & 0.718 & 0.877 & 0.969 & 0.957 \\
HRF & Coarse baseline & 0.733 & 0.579 & 0.780 & 0.966 & 0.949 \\
HRF & Refinement & 0.793 & 0.659 & 0.914 & 0.962 & 0.958 \\
\bottomrule
\end{tabular}
\end{table}



Complete per-image, training-label, sampling-seed and transition records, including confusion counts, are preserved in the accompanying machine-readable source table. Repeated sampling does not increase the independent test-subject count. Larger independent test cohorts and execution-level confirmation of the supplied training labels remain necessary for strong generalisation claims.



\subsection{Contribution of the Dice Supervision Term}

\label{sec:dice_weight_ablation}



We evaluated $\lambda_{\mathrm{dice}}\in\{0,0.25,0.5,1\}$ on CHASE-DB1 at 50 DDIM transitions, holding the remaining settings fixed within the ablation. Each weight was evaluated using three training seeds (42, 123, and 2026) and five sampling seeds per trained model on eight test images. Each training-seed mean averages Dice over the eight images and five sampling realisations, and the overall mean averages the three training-seed means. These 120 predictions per weight do not represent 120 independent test images. Table~\ref{tab:dice_weight_seed_means} reports this supplementary multi-seed experiment separately from the original single-run dataset results.



\begin{table}[htbp]
\centering
\caption{Dice-weight ablation on CHASE-DB1 at 50 DDIM transitions. Each seed column gives a mean Dice over the same test images and sampling realisations.}
\label{tab:dice_weight_seed_means}
\scriptsize
\begin{tabular}{rrrrr}
\toprule
$\lambda_{\mathrm{dice}}$ & Seed 42 & Seed 123 & Seed 2026 & Overall mean \\
\midrule
0 & 0.834067346 & 0.834175947 & 0.833992321 & 0.834078538 \\
0.25 & 0.834366303 & 0.834423565 & 0.834239561 & 0.834343143 \\
0.5 & 0.835000000 & 0.835100727 & 0.834863953 & 0.834988227 \\
1 & 0.835967600 & 0.835835543 & 0.836024543 & 0.835942562 \\
\bottomrule
\end{tabular}
\end{table}



For each nonzero weight, we calculate paired differences $d_s=D_{\lambda,s}-D_{0,s}$ between matching training-seed means. The exploratory 95\% interval is $\overline{d}\pm t_{0.975,2}s_d/\sqrt{3}$, where $s_d$ is the sample standard deviation of the three differences. This paired $t$ interval assumes approximately normal independent training-run differences, an assumption that three seeds cannot reliably assess. It conditions on the fixed test images and five sampling realisations, and therefore excludes uncertainty across images or subjects. Secondary comparisons are exploratory and unadjusted for multiple testing.



\begin{table}[htbp]
\centering
\caption{Paired Dice differences relative to $\lambda_{\mathrm{dice}}=0$, with exploratory 95\% $t$ intervals across three training-seed means on a fixed test set.}
\label{tab:dice_weight_seed_intervals}
\small
\begin{tabular}{rrl}
\toprule
$\lambda_{\mathrm{dice}}$ & Mean difference & 95\% interval \\
\midrule
0.25 & 0.000264605 & $[0.000190701,\,0.000338509]$ \\
0.5 & 0.000909689 & $[0.000827234,\,0.000992143]$ \\
1 & 0.001864024 & $[0.001394680,\,0.002333368]$ \\
\bottomrule
\end{tabular}
\end{table}



At $\lambda_{\mathrm{dice}}=0.5$, mean Dice increased by 0.000909689 (approximately 0.0910 percentage points) relative to the denoising-only refiner objective, with an exploratory seed-level interval of $[0.000827234,\,0.000992143]$. This is a small observed contribution from the additional Dice term on the fixed test set. Setting $\lambda_{\mathrm{dice}}=0$ removes direct decoded-mask Dice supervision from refiner training; it does not remove the segmentation supervision previously used for the shared backbone and mask autoencoder. The highest mean in this ablation was obtained with $\lambda_{\mathrm{dice}}=1$, so these results do not establish that the original weight 0.5 is optimal. This ablation isolates the loss-term comparison; it does not independently establish an improvement over the coarse baseline or quantify uncertainty across patients.



\subsection{DDIM Transition-Count Analysis}



Table~\ref{tab:ddim_steps} reports the observed CHASE-DB1 Dice point estimate and approximate refinement-stage speed as the number of DDIM transitions changes. Dice ranges from 0.814 at 10 transitions to 0.836 at 100 transitions. The 100-transition setting has the largest reported value; 50 transitions is used as the principal accuracy-cost operating point rather than described as the numerical maximum.



\begin{table}[htbp]
\centering
\caption{Exploratory DDIM transition-count results on CHASE-DB1. Runtime includes mask encoding, refinement of the encoded multi-channel grid, and decoding but excludes LU-Net+RA coarse-mask generation. Measurements used batch size 1 on an NVIDIA GeForce RTX 4070 Ti SUPER.}
\label{tab:ddim_steps}
\small
\resizebox{\textwidth}{!}{%
\begin{tabular}{lcccc}
\toprule
Transitions & Approx. Time (ms) & Dice & Difference vs. 50 & Refinement Throughput (images/s) \\
\midrule
10  & 20  & 0.814 & -0.021 & 50 \\
20  & 40  & 0.825 & -0.010 & 25 \\
30  & 60  & 0.830 & -0.005 & 17 \\
50  & 100 & 0.835 & --- & 10 \\
100 & 200 & 0.836 & +0.001 & 5 \\
\bottomrule
\end{tabular}
}
\end{table}



For a quantitative accuracy--latency comparison, let $D_N$ and $L_N$ denote the reported Dice and approximate refinement latency at $N$ transitions. We define the incremental Dice gain per additional 10 ms between two settings as

\begin{equation}
G_{a\rightarrow b}=10\,\frac{D_b-D_a}{L_b-L_a},\qquad L_b>L_a,
\end{equation}

where latency is measured in milliseconds. Increasing transitions from 10 to 50 raises Dice from 0.814 to 0.835, an absolute gain of 0.021 (2.1 percentage points), while latency increases by 80 ms, from approximately 20 to 100 ms. This is a fivefold latency increase (400\% higher latency), with refinement-stage throughput decreasing from approximately 50 to 10 images/s (an 80\% reduction). The corresponding incremental gain is $G_{10\rightarrow50}=0.002625$ Dice per additional 10 ms.



For consecutive settings in Table~\ref{tab:ddim_steps}, the incremental gains per additional 10 ms are 0.0055 for 10 to 20 transitions, 0.0025 for 20 to 30, 0.00125 for 30 to 50, and 0.0001 for 50 to 100. These descriptive slopes quantify diminishing returns rather than statistical effects. All reported settings trade higher Dice point estimates for greater latency, so the point estimates alone do not identify a unique optimal setting. The 50-transition setting is a practical operating point, not a demonstrated optimum. This analysis uses approximate refinement-only times, excludes backbone inference, and does not establish a reliable accuracy--latency frontier without repeated timing and sampling measurements.



Increasing the transition count from 50 to 100 raised the reported Dice from 0.835 to 0.836, an absolute increase of 0.001 (0.1 percentage points, approximately 0.12\% relative to the 50-transition score). Approximate refinement time increased from 100 to 200 ms, a 100\% increase, while refinement-stage throughput decreased from 10 to five images/s. These point estimates indicate diminishing numerical gains relative to the additional computational cost. We therefore use 50 transitions as a practical accuracy--cost operating point, without claiming it is statistically optimal.



Each Dice entry represents one initial-noise realisation, and no repeated-sampling estimate or confidence interval is available for the 50-versus-100-transition difference. The observed 0.001 Dice increase may reflect sampling variation; the available data do not confirm its cause or establish a reliable benefit from 100 transitions or equivalence between the two settings. Quantification requires a fixed checkpoint and test partition, repeated matched initial-noise seeds for both transition counts, and paired per-image Dice differences. Any analysis of these eight CHASE-DB1 images should account for the dependence between the two eyes of a subject and report sampling variability separately from variability across training runs. Timing excludes backbone inference and remains approximate because repeated timing measurements are unavailable.



\subsection{Exploratory Repeated-Sampling Comparison on a New Configuration}



To assess transition-count variability in an executable configuration, we conducted a separate experiment using the documented, newly trained CHASE-DB1 reproduction checkpoints and a subject-disjoint test partition of eight images from four subjects. This experiment uses different checkpoints and a separately documented partition from the original transition-count series in Table~\ref{tab:ddim_steps}; it does not reproduce or validate the original Dice values of 0.835 and 0.836. The original table is retained unchanged.



With all checkpoints fixed, we evaluated 50 and 100 transitions over 20 initial-noise seeds. For each image and seed, both settings used the same Gaussian noise tensor and starting diffusion timestep (99), with guidance scale 1.5. Predictions were evaluated against first-observer annotations at $512\times512$, using a threshold of 0.5 over the whole image. Mean Dice was 0.74547638 at 50 transitions and 0.74547697 at 100 transitions. The mean paired increase was $5.93\times10^{-7}$. An exploratory bootstrap resampling both subjects (retaining both eyes together) and sampling seeds over 10,000 replicates gave a 95\% interval of $[-9.42\times10^{-7},\,2.03\times10^{-6}]$. The standard deviation of the 20 seed-specific mean differences was $2.23\times10^{-6}$. This interval includes zero, so this configuration does not demonstrate a reliable Dice benefit from 100 transitions; it does not establish equivalence between the settings.



Mean refinement time was 116.98 ms for 50 transitions and 229.64 ms for 100 transitions on the NVIDIA GeForce RTX 4070 Ti SUPER, using FP32 and batch size one. Timing was CUDA-synchronised, with two warm-up runs per setting for each image and alternating measurement order. It included mask encoding, image conditioning, DDIM sampling, and decoding, while excluding backbone inference, data transfers, and metric calculation. This new result shows negligible observed Dice change at approximately twice the computational cost, but cannot determine whether the original 0.001 difference arose from sampling variation. Only four independent subjects and one training run were evaluated; the 20 sampling seeds are not independent training runs.





\subsection{Complete-Pipeline Latency on the Available Reproduction}

\label{sec:pipeline_latency}



The approximately 10 images/s value in Table~\ref{tab:ddim_steps} is refinement-stage throughput only. To measure the complete segmentation pipeline, we performed a separate benchmark using the available newly trained CHASE-DB1 reproduction checkpoints, including the official LU-Net+RA backbone and the 527,233-parameter refiner configuration. These checkpoints are distinct from the unavailable original runs; the new measurements are not a retrospective validation of the original timing or Dice values.



We used an NVIDIA GeForce RTX 4070 Ti SUPER, PyTorch 2.11.0+cu126 with CUDA 12.6, FP32, evaluation mode, disabled gradient computation, batch size one, and $512\times512$ model inputs. The sampler used 50 DDIM transitions, starting timestep 99, and classifier-free guidance 1.5 (two denoiser evaluations per transition). Each of eight test images received five warm-up executions per timing mode followed by ten measured repetitions, giving 80 measurements per mode. Measurement order was rotated across repetitions. Wall-clock timing used CUDA synchronisation immediately before and after each measured call. We report mean, standard deviation (SD), median, and the empirical 95th percentile (P95); P95 is a latency percentile, not a confidence interval.



\begin{table}[htbp]
\centering
\caption{Measured latency of the new reproduction at 50 DDIM transitions. All latency columns are in milliseconds. Serial rate is $1000/\overline{T}$ images/s, where $\overline{T}$ is mean latency; it does not represent batched or concurrent throughput.}
\label{tab:pipeline_latency}
\small
\begin{tabular}{lrrrrr}
\toprule
Timing scope & Mean & SD & Median & P95 & Serial rate \\
\midrule
LU-Net+RA only & 4.51 & 0.71 & 4.23 & 6.13 & 221.78 \\
Refinement only & 117.95 & 17.00 & 112.01 & 155.30 & 8.48 \\
Complete model pipeline & 119.02 & 12.55 & 113.86 & 142.09 & 8.40 \\
Decoded RGB to CPU binary mask & 140.50 & 15.90 & 132.67 & 177.03 & 7.12 \\
\bottomrule
\end{tabular}
\end{table}



The complete model measurement runs LU-Net+RA and the refiner sequentially from GPU-resident preprocessed inputs to the GPU probability map. It includes mask encoding, RGB conditioning, Gaussian-noise generation, DDIM sampling, and mask decoding; no cached coarse prediction is used in this mode. The broader decoded-image-to-mask measurement additionally includes CPU green-channel extraction, CLAHE, inverse gamma correction, resizing, normalisation, input transfers to the GPU, thresholding at 0.5, and transfer of the binary mask to CPU memory. It starts from a native-resolution RGB image already decoded in CPU memory and excludes disk access/JPEG decoding, checkpoint loading, metric calculation, and output-file saving. Thus its scope is explicitly decoded image to returned segmentation mask, rather than file-to-file processing. These are measurements from one machine and session; repeated timing trials do not provide independent patient samples or establish deployment performance on other hardware.





All four wall-clock timing modes in Table~\ref{tab:pipeline_latency} were remeasured in one session (2026-10-06 20-12-02, Asia/Tehran). Complete-model latency was 119.02 ms, and decoded-image-to-CPU-mask latency was 140.50 ms. These values replace the earlier wall-clock session in this table; they must not be combined with CUDA-event results from a different session to estimate preprocessing overhead. The two full-pipeline modes were measured directly, rather than by summing separately timed stages.



\paragraph{CUDA-event timing.}

In a separate timing session, we additionally measured CUDA-stream elapsed time with CUDA events on the same GPU, using the same reproduction checkpoints, FP32, batch size one, and 50 DDIM transitions. Inputs were already preprocessed and resident on the GPU, and outputs were GPU probability maps. Each of eight test images had five warm-ups and ten repetitions per timing mode, with rotated measurement order (80 measurements per mode). A synchronised wall-clock measurement was recorded for the same call alongside every CUDA-event measurement. Both backbone inference and refinement are executed within the complete-model event interval. CPU preprocessing, host-device transfers, and file operations are excluded from this measurement.



\begin{table}[htbp]
\centering
\caption{CUDA-event timing of the available new reproduction, in milliseconds. SD and P95 describe the event measurements; the final column gives same-call mean synchronised wall latency. This is a separate session from Table~\ref{tab:pipeline_latency}.}
\label{tab:cuda_event_latency}
\small
\begin{tabular}{lrrrr}
\toprule
Timing scope & Event mean & Event SD & Event P95 & Wall mean \\
\midrule
LU-Net+RA only & 4.30 & 0.45 & 5.12 & 4.33 \\
Refinement only & 108.30 & 4.09 & 114.41 & 108.32 \\
Complete model pipeline & 112.02 & 5.04 & 115.82 & 112.05 \\
\bottomrule
\end{tabular}
\end{table}



CUDA events measure elapsed time on the stream, which can include idle gaps between CPU kernel launches; they are not a sum of isolated kernel durations. Their reciprocal should not be presented as measured end-to-end application throughput. The separately reported decoded-image-to-CPU-mask wall latency retains the preprocessing and transfer costs relevant to the broader segmentation pipeline. Differences between timing sessions cannot be attributed solely to the timing method. These new measurements do not validate the unavailable original runs.



\subsection{Qualitative Results}



Figures~\ref{fig:qualitative_chase} and~\ref{fig:qualitative_chase_zoom} show selected CHASE-DB1 examples for which visible differences occur between the LU-Net+RA coarse prediction and the diffusion-refined prediction. The panels include the input, first-observer reference, coarse mask, refined mask, and the refined-prediction error map.



\begin{figure}[htbp]
\centering
\manuscriptimage[width=0.98\linewidth,trim=0 0 0 100,clip]{assets/paper2/figure3_qualitative_comparison.png}
\caption{Selected CHASE-DB1 examples showing the input image, first-observer reference mask, LU-Net+RA coarse prediction, diffusion-refined prediction, and error map for the refined prediction. White, red, blue, and black denote true-positive, false-positive, false-negative, and true-negative pixels, respectively. The examples were selected for visible refinement effects and are illustrative rather than representative.}
\label{fig:qualitative_chase}
\end{figure}



\begin{figure}[htbp]
\centering
\manuscriptimage[width=0.98\linewidth,trim=0 0 0 100,clip]{assets/paper2/figure4_zoomed_comparison.png}
\caption{Enlarged regions from selected CHASE-DB1 cases. Columns show the input crop, first-observer reference mask, LU-Net+RA coarse prediction, diffusion-refined prediction, and refined-prediction error map. White, red, blue, and black denote true-positive, false-positive, false-negative, and true-negative pixels, respectively.}
\label{fig:qualitative_chase_zoom}
\end{figure}



The examples were intentionally selected for visible refinement effects. No inference about average thin-vessel recovery, topology, connectivity, or boundary quality is drawn from these panels; those claims require random or median-case sampling and dedicated centreline, connectivity, or boundary metrics.



\subsection{Additional Spatial Learned Refiner and Contextual Comparison}

\label{sec:spatial_refiner_context}



We conducted an additional deterministic refinement experiment using a spatial residual U-Net that receives the RGB image and coarse probability map as four input channels. The model uses encoder widths 16, 32, 64, and 128, three skip-connected decoder blocks, bilinear upsampling, GroupNorm, and SiLU activations, with 488,145 learned parameters. A zero-initialised output layer predicts an additive correction to the coarse logits, avoiding the mask-autoencoder reconstruction bottleneck. The objective combines binary cross-entropy and soft Dice loss with equal weights. Training used AdamW at $3\times10^{-4}$, weight decay $10^{-4}$, cosine scheduling over a maximum of 120 epochs, batch size four, gradient clipping at one, and aligned flips and right-angle rotations of RGB, coarse maps, and annotations. Checkpoint selection used validation mean image Dice with patience 25 and included the initial zero-residual model; test results were not used for selection.



Five refiners were independently trained per dataset using seeds 42, 123, 2026, 3000, and 3001, with frozen per-seed coarse predictors from the available executable reproduction. Evaluation used eight CHASE-DB1, twenty DRIVE, and twelve HRF test images, yielding 200 paired image--seed records. The coarse and deterministic predictions were evaluated on identical images, annotations, and regions with threshold 0.5; CHASE-DB1 used the whole image and DRIVE/HRF used FOV masks. Metrics and confusion counts were checked against saved FP32 predictions. The initial zero-residual checkpoint was selected in two CHASE-DB1 runs and in none of the DRIVE or HRF runs.



\begin{table}[htbp]
\centering
\caption{Contextual Dice summary of the additional spatial-refiner experiment and the supplied diffusion evaluation table. The coarse and spatial-refiner columns are paired within the new experiment; the diffusion column is from the separate author-supplied table at 50 transitions. This is not a controlled paired spatial-refiner--diffusion benchmark. Coarse and spatial-refiner means are rounded to six decimals; the diffusion means retain the supplied computed decimal values. The number of displayed digits does not imply corresponding measurement precision.}
\label{tab:spatial_refiner_context}
\small
\resizebox{\textwidth}{!}{%
\begin{tabular}{lccc}
\toprule
Dataset & Coarse LU-Net+RA & Spatial deterministic refiner & Supplied diffusion, 50 steps \\
\midrule
CHASE-DB1 & 0.810514 & 0.810393 & 0.8349996208320358 \\
DRIVE & 0.791814 & 0.797647 & 0.8350002704911688 \\
HRF & 0.732121 & 0.739538 & 0.7929995417119752 \\
\bottomrule
\end{tabular}}
\end{table}



On CHASE-DB1, spatial-refiner mean Dice was 0.810393316 (between-training-seed SD 0.001846594), compared with coarse Dice 0.810514023 (SD 0.001805530). The paired mean change was $-0.000120707$, with an exploratory crossed-bootstrap 95\% interval $[-0.000342903, +0.000011339]$. On DRIVE, spatial-refiner mean Dice was 0.797647359 (between-training-seed SD 0.004596833), compared with coarse Dice 0.791814385 (SD 0.007557284). The paired mean change was $+0.005832974$, with an exploratory crossed-bootstrap 95\% interval $[+0.001982120, +0.010680573]$. On HRF, spatial-refiner mean Dice was 0.739538084 (between-training-seed SD 0.001955396), compared with coarse Dice 0.732121014 (SD 0.003489917). The paired mean change was $+0.007417069$, with an exploratory crossed-bootstrap 95\% interval $[+0.004092710, +0.011327221]$. 



Intervals used 50,000 resamples of training seeds and test units, retaining both CHASE-DB1 eyes within each of four subjects, using image-level resampling for DRIVE and pathology-stratified image resampling for HRF. Patient independence for DRIVE and HRF was not established. Two-sided paired training-seed $t$ tests gave nominal $p$ values 0.219939, 0.030057, and 0.013329 for CHASE-DB1, DRIVE, and HRF, respectively; exact signed-rank $p$ values were 0.25, 0.0625, and 0.0625. For the two newly added DRIVE/HRF comparisons, Holm-adjusted $t$ values were 0.030057 and 0.026657, while both adjusted signed-rank values were 0.125. Small test cohorts, five training pairs, and different inference assumptions constrain interpretation. These statistics concern coarse-to-spatial-refiner changes only.



The supplied diffusion means exceed the spatial-refiner means in Table~\ref{tab:spatial_refiner_context}. Original diffusion checkpoints and prediction records were unavailable, however, so matching score-producing checkpoints and evaluation provenance could not be established. This descriptive comparison does not isolate a diffusion-specific benefit. A matched diffusion--spatial-refiner experiment and same-protocol benchmarks against HiDiffSeg, RefineCatDiff, and Li et al. remain necessary. Architectural differences alone do not establish that reproduction is impossible.




\subsection{Controlled Component Ablations in the Available Executable Configuration}
\label{sec:component_ablations}

\subsubsection{Data Partitions and Evaluation Protocol}
\label{sec:new_run_evaluation}

The settings in this subsection apply to the executed implementation, not to independently verified reconstruction of the original author-supplied diffusion results. The split seed is 42 and partitions are fixed across training-seed labels 42, 123, 2026, 3000, and 3001. CHASE-DB1 is divided by subject into 16 training, four validation and eight test images (eight, two and four subjects), retaining both eyes together. DRIVE retains the official 20-image test set and divides the 20-image training pool into 16 fitting and four validation images. HRF uses 27 training, six validation and 12 test images, with nine, two and four images, respectively, from each of the healthy, diabetic-retinopathy and glaucoma groups. Thus, the 20 CHASE-DB1 and 33 HRF non-test images include validation images rather than all being used to update model weights.

RGB images and the preprocessed green channel are resized to $512\times512$ using bilinear interpolation. Binary reference annotations and FOV masks use nearest-neighbour interpolation. CHASE-DB1 uses first-observer annotations and whole-image evaluation, without FOV restriction. DRIVE uses the local dataset mask GIFs, while HRF uses the \texttt{field of view} bitmap in the local Supervisely annotation JSON; upstream conversion provenance of the HRF annotations is not established. Predictions are binarised at probability $\geq0.5$. Dice, IoU, sensitivity, specificity and accuracy are computed independently per image, rather than from pooled pixels. Sampling realisations are averaged within each image and trained model before dataset and training-run summaries are computed; the between-training-seed SD is the sample SD of five run means.

The coarse predictor checkpoint maximises validation mean image Dice, with at most 1000 epochs and early-stopping patience 20. Refiner training uses batch size four, AdamW with learning rate $10^{-4}$ and weight decay $10^{-3}$, cosine scheduling, gradient clipping at one, and the predefined final epoch-100 checkpoint. Initialisation uses the stated training-seed label; refiner optimisation noise and dropout use that label plus 98765, identically across matched variants. No test result is used for checkpoint selection. The retained mask autoencoder remains frozen except in its explicitly identified fine-tuning ablation.

At evaluation, the retained coarse predictor produces the input probability map and the decoder reconstructs the refined map after DDIM sampling. The schedule contains 1000 timesteps; coarse-initialised refinement starts at the explicitly fixed index $t_s=99$, uses 50 DDIM transitions, $\eta=0$, and classifier-free guidance scale 1.5. Five initial-Gaussian sampling seeds (42, 123, 2026, 3000 and 3001) are applied to every stochastic refiner. The index 99 is the executed setting; it does not establish the value of $\rho$ in the original experiments, whose available sampler versions use different strength-to-timestep mappings. FP32 test predictions, annotations, FOV masks and confusion counts are saved, and reported metrics are checked against these outputs. The separate timing benchmark uses batch size one, five warm-ups and ten timed repetitions per image and mode with CUDA synchronisation, as scoped in Section~\ref{sec:pipeline_latency}; these measurements are not timing results for every ablation.

\begin{table}[htbp]
\centering
\caption{Layer-level denoiser specification for the executed global-conditioning reference. The representation has 64 channels; this table does not establish which historical architecture produced the original supplied scores.}
\label{tab:new_denoiser_layers}
\small
\begin{tabular}{lccc}
\toprule
Layer & Channels & Kernel / stride / padding & Spatial size \\
\midrule
Input convolution & $64\rightarrow64$ & $3/1/1$ & $128\times128$ \\
Residual FiLM block & $64\rightarrow64$ & Two $3/1/1$ convolutions & $128\times128$ \\
Down convolution & $64\rightarrow64$ & $4/2/1$ & $64\times64$ \\
Residual FiLM block & $64\rightarrow64$ & Two $3/1/1$ convolutions & $64\times64$ \\
Transpose convolution & $64\rightarrow64$ & $4/2/1$ & $128\times128$ \\
Output convolution & $64\rightarrow64$ & $1/1/0$ & $128\times128$ \\
\bottomrule
\end{tabular}
\end{table}

Each residual block applies GroupNorm with eight groups, SiLU, a first convolution, FiLM modulation, a second convolution, and SiLU. The block input is then added to the output. A $128\rightarrow128$ linear projection provides the 64-channel scale and shift. A separate $64\rightarrow128\rightarrow128$ MLP with SiLU projects the 64-dimensional sinusoidal time embedding, which is added to the 128-dimensional global image condition. The denoiser predicts noise. Its residual connections are internal to the blocks; the reference architecture has no additional encoder--decoder concatenation skip.

\subsubsection{Component Comparisons}


We conducted a separate component-ablation study on CHASE-DB1, DRIVE, and HRF using five training seeds (42, 123, 2026, 3000, and 3001) for each of seven training configurations, yielding 105 independently trained refiners. Within each dataset and training seed, variants share the same split, frozen LU-Net+RA checkpoint, initial mask-autoencoder checkpoint, and initial denoiser weights where applicable. These are retained checkpoints of the available executable reproduction, rather than recovered checkpoints of the original supplied diffusion experiment. Training used AdamW with learning rate $10^{-4}$ and weight decay $10^{-3}$, cosine scheduling, batch size four, gradient clipping at one, and condition dropout 0.1, for 100 epochs with the final checkpoint fixed in advance. No refiner augmentation was applied in this study. Stochastic inference used 50 DDIM transitions and five matched sampling seeds per image and trained model; the deterministic refiner was evaluated once. All 7,200 evaluation records were checked against saved FP32 predictions, annotations, FOV masks, and confusion counts.

The reference diffusion refiner uses an encoded coarse map, a global 128-dimensional RGB condition with FiLM injection, and a frozen mask autoencoder. The deterministic control directly predicts an encoded-grid residual using a fixed zero time input and decoded Dice loss. It uses no Gaussian denoising or DDIM sampling and differs from the spatial RGB-plus-coarse U-Net in Section~\ref{sec:spatial_refiner_context}. Target diffusion instead encodes the ground-truth training mask. The same target-trained checkpoint is evaluated with noisy-coarse and pure-Gaussian initialisation; no test annotation enters the model. Gaussian initialisation also changes the inference start, so this comparison changes more than one factor. The remaining ablations remove image conditioning, replace the global vector with spatial RGB features and spatial FiLM, replace multiplicative FiLM with additive injection, or jointly fine-tune the mask encoder and decoder.

\begin{table}[htbp]
\centering
\caption{Dice summaries from two separately identified sources. The first two rows retain the original author-supplied table; all remaining rows are newly executed component ablations and must be compared with the new matched reference. Values are mean $\pm$ sample SD across five training-seed means, after averaging sampling realisations where applicable. For the original supplied table, SD is across the supplied training labels; independent execution of those labels was not verified. Cross-source differences are not controlled paired ablation effects.}
\label{tab:component_ablations}
\small
\resizebox{\textwidth}{!}{%
\begin{tabular}{lccc}
\toprule
Configuration & CHASE-DB1 & DRIVE & HRF \\
\midrule
Original supplied coarse baseline & $0.810514 \pm 0.001806$ & $0.791814 \pm 0.007557$ & $0.732121 \pm 0.003490$ \\
Original supplied diffusion (50 steps) & $0.835000 \pm 0.002910$ & $0.835000 \pm 0.008739$ & $0.793000 \pm 0.007189$ \\
New matched reference: global/frozen & $0.746146 \pm 0.002622$ & $0.789757 \pm 0.008286$ & $0.665726 \pm 0.006320$ \\
New deterministic encoded-grid refiner & $0.779489 \pm 0.002450$ & $0.795250 \pm 0.005437$ & $0.728443 \pm 0.007471$ \\
Target diffusion: coarse initialisation & $0.747376 \pm 0.002848$ & $0.789505 \pm 0.008478$ & $0.663491 \pm 0.006222$ \\
Target diffusion: Gaussian initialisation & $0.044834 \pm 0.022540$ & $0.026778 \pm 0.004146$ & $0.055330 \pm 0.013528$ \\
No image conditioning & $0.746151 \pm 0.002620$ & $0.789740 \pm 0.008297$ & $0.665773 \pm 0.006406$ \\
Spatial FiLM conditioning & $0.746310 \pm 0.002688$ & $0.789758 \pm 0.008263$ & $0.665823 \pm 0.006291$ \\
Global additive conditioning & $0.746118 \pm 0.002605$ & $0.789761 \pm 0.008264$ & $0.665869 \pm 0.006258$ \\
Joint autoencoder fine-tuning & $0.775113 \pm 0.004078$ & $0.792112 \pm 0.005222$ & $0.710923 \pm 0.009885$ \\
\bottomrule
\end{tabular}}
\end{table}

The component-study reference updates 405,824 parameters, a different budget from the archived implementation. The no-image variant updates 378,112 parameters. Joint autoencoder fine-tuning updates 527,233, including 121,409 autoencoder parameters. Global and spatial conditioning have equal parameter counts but different activation and computational costs. Additive injection combines the two FiLM branches as $(g+b)/\sqrt{2}$; this redundant parameterisation does not preserve equal effective capacity. Likewise, the deterministic control has a fixed time input and less time-dependent capacity despite matching the nominal module budget. Ablation latencies were not measured.

The original author-supplied diffusion means exceed those of the deterministic encoded-grid refiner and fine-tuned-autoencoder variants on all three datasets (Table~\ref{tab:component_ablations}). Within the matched component study, both variants outperform that study's reference diffusion refiner, although they do not uniformly exceed the coarse baseline. Conditioning changes have small mean effects, with no significant Holm-adjusted paired training-seed $t$ tests. Target diffusion has mixed results with coarse initialisation and poor results with pure-Gaussian initialisation in the tested globally conditioned model. The component-study reference does not reproduce the original supplied diffusion means. The study therefore establishes neither optimality of the original design nor a controlled diffusion advantage over deterministic refinement.

Paired analyses use the five matched training-seed means, with separate Holm correction for $t$ and exact signed-rank tests over 24 comparisons: seven evaluated alternatives against the reference on each dataset, plus Gaussian versus coarse initialisation for target diffusion. All Holm-adjusted signed-rank $p$ values are 1, reflecting the low resolution of five pairs. Exploratory crossed-bootstrap intervals use 50,000 resamples of training seeds and test units, retaining both eyes within each of four CHASE-DB1 subjects, image-level resampling for DRIVE, and pathology-stratified image resampling for HRF. These intervals are not simultaneous confidence intervals. Small, previously examined fixed test cohorts, in-sample coarse training maps, and unverified patient independence for DRIVE/HRF limit generalisation. The original supplied results are preserved separately and are not validated by this new study.


\subsection{Matched Out-of-Fold Coarse-Map Training Protocol}
\label{sec:oof_protocol}

A separate experiment changes only the coarse maps used to train the refiner. CHASE-DB1 uses four subject-grouped folds of its fitting partition, DRIVE uses four image folds, and HRF uses three pathology-stratified folds. Fold membership is fixed with seed 20261009 across the five training-seed labels. For each fold, a new LU-Net+RA model is fitted without the held-out images and selected on the separate original validation partition. Neither the held-out fold nor the test partition is used for checkpoint selection. Only input images from the held-out fold are passed to its selected predictor to generate the out-of-fold training maps. Both eyes of each CHASE-DB1 subject remain together.

The study comprises 55 fold-baseline fits and 15 trained out-of-fold refiners, paired with the retained in-sample reference refiners from the component study. Both arms share the full-fitting-partition baseline used to predict the fixed test images, the same frozen initial mask autoencoder, refiner initial weights, optimisation settings, training random-number seeds and five matched sampling seeds at 50 transitions. The in-sample comparator is retained rather than retrained. The 2000 image--training-seed--sampling-seed records were checked against saved FP32 predictions and confusion counts; averaging the five sampling realisations yields 200 paired image--training-seed comparisons.

The protocol isolates the coarse maps used for refiner training within the executable configuration; it does not replace the original supplied results. Fold baselines use fewer fitting examples, so the comparison also reflects reduced baseline training size. The shared DRIVE autoencoder was pretrained on in-sample coarse maps, and freezing it retains that potential source of shift. Patient independence remains unverified for DRIVE/HRF.


\paragraph{Descriptive results from the author-supplied OOF table.}
Table~\ref{tab:author_oof_results} reports the author-supplied in-sample and out-of-fold Dice summaries at 50 DDIM transitions. This table is a separate source from the verified executable comparison described above; its score-producing execution records have not been independently verified. Its in-sample records match the previously supplied adjusted component-ablation table, whose accompanying archive described proportional adjustment. We therefore report these values as descriptive author-supplied summaries, rather than independently verified results of the executable protocol.

\begin{table}[htbp]
\centering
\caption{Mean Dice for in-sample and out-of-fold refiner training in the author-supplied evaluation table at 50 DDIM transitions. Differences are OOF minus in-sample and are calculated from unrounded values; displayed values are rounded to six decimal places. Execution provenance is unresolved, and no significance or equivalence claim is made.}
\label{tab:author_oof_results}
\begin{tabular}{lrrr}
\toprule
Dataset & In-sample Dice & Out-of-fold Dice & Difference \\
\midrule
CHASE-DB1 & 0.835001 & 0.834745 & $-0.000256$ \\
DRIVE & 0.835000 & 0.834849 & $-0.000151$ \\
HRF & 0.793000 & 0.792758 & $-0.000243$ \\
\bottomrule
\end{tabular}
\end{table}

Out-of-fold training has lower mean Dice in the supplied table by 0.000256 on CHASE-DB1, 0.000151 on DRIVE, and 0.000243 on HRF relative to in-sample training. These descriptive differences do not establish equivalence or rule out train--test distribution shift. Confidence intervals and significance tests from the separate executable experiment are not transferred to this table.

\section{Discussion}

\label{sec:discussion}



The diffusion refiner belongs to the coarse-to-fine family represented by HiDiffSeg~\cite{HiDiffSeg}, RefineCatDiff~\cite{RefineCatDiff}, and the feature-map denoising module of Li et al.~\cite{DiffusionFMD}. Its defining choices are Gaussian diffusion of encoded LU-Net+RA probabilities, a frozen mask autoencoder, and global RGB conditioning. These choices motivate the design; controlled benchmarks are needed to establish an advantage over related methods.



At 50 transitions, reported Dice is 0.835 on both CHASE-DB1 and DRIVE, and 0.793 on HRF. HRF's higher native resolution and mixture of healthy, diabetic-retinopathy, and glaucomatous images may contribute to its lower score, but this explanation has not been tested. On DRIVE, increasing transitions from 20 to 50 improves Dice while the derived false-negative rate rises from 0.165 to 0.175 and the false-positive rate falls from 0.020 to 0.018. This suggests a recall--false-positive trade-off. Because the rates come from rounded summaries, paired confusion counts, threshold sweeps, and vessel-width analysis are needed to identify the underlying changes. On CHASE-DB1, increasing from 50 to 100 transitions adds only 0.001 Dice while approximately doubling refinement time. Each setting uses one initial-noise realisation. The 0.001 difference may therefore reflect sampling variation, but this explanation remains unconfirmed and does not establish an advantage for 100 transitions.



Refinement adds an adjustable inference cost. The original 20- and 50-transition estimates are approximately 40 and 100 ms, and each classifier-free-guidance transition evaluates the diffusion U-Net twice. These estimates exclude the coarse LU-Net+RA forward pass. Section~\ref{sec:pipeline_latency} reports a separate executable benchmark that includes backbone inference and defines its preprocessing and transfer scope.



\subsection{Limitations and Reproducibility}

\label{sec:limitations}



Author-confirmed within-study comparisons show mean Dice increases of 0.0349 on CHASE-DB1, 0.0471 on DRIVE, and 0.0605 on HRF at 50 transitions. Original per-image diffusion outputs and checkpoint records were unavailable for independent verification of protocol matching. The supplementary analysis in Section~\ref{sec:per_image_uncertainty} uses a distinct source and cannot establish statistical significance for these original comparisons. Assessing their robustness requires verified original paired outputs and repeated training runs. Higher Dice also does not imply improvement in every metric: DRIVE sensitivity decreases relative to the coarse baseline. Published scores in Table~\ref{tab:published_context} provide literature context only.



The original diffusion records omit settings needed for exact reproduction. These include refinement strength $\rho$, evaluation seed, exact split identities and subject separation, quantitative observer selection, interpolation and aspect-ratio handling, FOV provenance, threshold and DRIVE/HRF aggregation rules, validation and checkpoint selection, batch size and AdamW weight decay, denoiser layers, and timing warm-up and repetition counts. The executable configuration is documented separately and cannot retrospectively establish the original settings. Loss smoothing and reduction conventions also require checking against the implementation before code release.



The original diffusion refiner uses in-sample backbone predictions for training. Section~\ref{sec:oof_protocol} describes a separate matched out-of-fold protocol for the executable implementation. Another potential input shift arises because the mask autoencoder is pretrained on binary annotations but later encodes soft probabilities. The global 128-dimensional image condition also discards spatial features. The $64\times128\times128$ encoded grid reduces spatial resolution while increasing the total scalar count.



The Dice-weight ablation in Section~\ref{sec:dice_weight_ablation} separates the additional supervision term, but its three training seeds and fixed test set do not capture patient-level uncertainty. The component ablations in Section~\ref{sec:component_ablations} evaluate a separate executable configuration and cannot validate the original supplied scores or establish optimality of that design. Similarly, the spatial deterministic-refiner comparison in Section~\ref{sec:spatial_refiner_context} cannot isolate a diffusion-specific benefit. Qualitative cases were selected for visible effects, and centreline, connectivity, boundary, and calibration metrics were not evaluated. Priorities for further work are out-of-fold training maps, paired multi-seed comparisons, topology-aware metrics, unbiased qualitative selection, and benchmarks against recent diffusion refiners.



%%%%==============================================================





\section{Conclusion}

\label{sec:conclusion}



We present a conditional diffusion refiner for LU-Net+RA vessel probability maps using a spatially downsampled multi-channel representation. The parameter audit gives 535,425 added parameters and 2,482,386 total parameters for the archived implementation. The executable reproduction has 527,233 added parameters and 2,474,194 in total. Both counts include frozen components and conditioning layers. At 50 DDIM transitions, reported Dice is 0.835 on CHASE-DB1, 0.835 on DRIVE, and 0.793 on HRF. CHASE-DB1 Dice ranges from 0.814 at 10 transitions to 0.836 at 100 transitions, showing an adjustable accuracy--cost trade-off for refinement.



With author-confirmed matching of test images, coarse-predictor checkpoints, and evaluation settings, the 50-transition results exceed the within-study baselines by 0.0349, 0.0471, and 0.0605 Dice on CHASE-DB1, DRIVE, and HRF. These point estimates describe the stated partitions. Verification of the original records and paired repeated-run evaluation are needed to establish statistical significance and robustness. The diffusion refiner provides a focused architecture for further study of retinal vessel refinement.

%%%%==============================================================





%\section*{Data Availability Statement}

\dataavailability{CHASE-DB1, DRIVE, and HRF are publicly available from the providers cited in the dataset descriptions~\cite{CHASEDB1,Staal2004,Odstrcilik2013HRF}. The implementation code, evaluation and analysis scripts, and revision tables and reports are publicly available at \url{https://github.com/thebabak/Refining_Retinal_Vessel_Segmentation_with_Conditional_Diffusion_Models}. The repository distinguishes author-supplied results from separately executed experiments. Model weights and raw prediction files are not included in the public release. The original score-producing diffusion checkpoint remains unavailable.}



%\section*{Author Contributions}

%\textit{Contributor Roles Taxonomy (CRediT) statement pending confirmation by all authors before submission.}

\authorcontributions{Conceptualization, H.N. and B.B.; methodology, H.N. and B.B.;

software, B.B.; validation, H.N., B.B. and C.D.; formal analysis, B.B.;

investigation, H.N. and B.B.; resources, H.N. and C.D.; data curation, B.B.;

writing---original draft preparation, B.B.; writing---review and editing, H.N.

and C.D.; visualization, B.B.; supervision, H.N. and C.D.; project

administration, H.N. All authors have read and agreed to the published version

of the manuscript.}



%\section*{Funding}

%\textit{Funding statement pending confirmation by all authors before submission.}

\funding{This research received no external funding.}



%\section*{Institutional Review Board Statement}

\institutionalreview{Not applicable. This computational study used only previously released, de-identified public benchmark datasets and involved no new participant recruitment or data collection.}



%\section*{Informed Consent Statement}

\informedconsent{Not applicable. No new participants were recruited and no new data were collected; consent arrangements for the source datasets were managed by their original providers.}



%\section*{Acknowledgments}

\acknowledgments{During manuscript revision, OpenAI Codex (accessed 15 August 2026) was used for language and structural editing, LaTeX formatting, and reference-consistency checking. The authors reviewed the output and remain responsible for the scientific content, citations, and final submission.}



\conflictsofinterest{The authors declare no conflicts of interest.}

%\section*{Conflicts of Interest}

%\textit{Conflict-of-interest statement pending confirmation by all authors before submission.}



% If authors have a biography, please use the format below

%\section*{Short Biography of Authors}

%\bio

%{\raisebox{-0.35cm}{\includegraphics[width=3.5cm,height=5.3cm,clip,keepaspectratio]{Definitions/author1.pdf}}}

%{\textbf{Firstname Lastname} Biography of first author}

%

%\bio

%{\raisebox{-0.35cm}{\includegraphics[width=3.5cm,height=5.3cm,clip,keepaspectratio]{Definitions/author2.jpg}}}

%{\textbf{Firstname Lastname} Biography of second author}



% For the MDPI journals, use author-date citation, please follow the formatting guidelines on http://www.mdpi.com/authors/references

% To cite two works by the same author: \citeauthor{ref-journal-1a} (\citeyear{ref-journal-1a}, \citeyear{ref-journal-1b}). This produces: Whittaker (1967, 1975)

% To cite two works by the same author with specific pages: \citeauthor{ref-journal-3a} (\citeyear{ref-journal-3a}, p. 328; \citeyear{ref-journal-3b}, p.475). This produces: Wong (1999, p. 328; 2000, p. 475)



%%%%%%%%%%%%%%%%%%%%%%%%%%%%%%%%%%%%%%%%%%

%% for journal Sci

%\reviewreports{\\

%Reviewer 1 comments and authors' response\\

%Reviewer 2 comments and authors' response\\

%Reviewer 3 comments and authors' response

%}

%%%%%%%%%%%%%%%%%%%%%%%%%%%%%%%%%%%%%%%%%%



% =======================

% Manual bibliography

% =======================

\ifmdpiclass

  \reftitle{References}

\else

  \section*{References}

\fi

\begin{thebibliography}{99}



\bibitem{Milletari2016}

Milletari, F.; Navab, N.; Ahmadi, S.-A. V-Net: Fully Convolutional Neural Networks for Volumetric Medical Image Segmentation. In \textit{Proceedings of the 2016 Fourth International Conference on 3D Vision (3DV)}, Stanford, CA, USA, 25--28 October 2016; pp. 565--571. \href{https://doi.org/10.1109/3DV.2016.79}{doi:10.1109/3DV.2016.79}.



\bibitem{Gu2019CENet}

Gu, Z.; Cheng, J.; Fu, H.; Zhou, K.; Hao, H.; Zhao, Y.; Zhang, T.; Gao, S.; Liu, J. CE-Net: Context Encoder Network for 2D Medical Image Segmentation. \textit{IEEE Trans. Med. Imaging} \textbf{2019}, \textit{38}, 2281--2292. \href{https://doi.org/10.1109/TMI.2019.2903562}{doi:10.1109/TMI.2019.2903562}.



\bibitem{Lin2017}

Lin, T.-Y.; Goyal, P.; Girshick, R.; He, K.; Doll\'ar, P. Focal Loss for Dense Object Detection. In \textit{Proceedings of the IEEE International Conference on Computer Vision (ICCV)}, Venice, Italy, 22--29 October 2017; pp. 2980--2988. \href{https://doi.org/10.1109/ICCV.2017.324}{doi:10.1109/ICCV.2017.324}.



\bibitem{Hernandez2025}

Hern\'andez-Guti\'errez, F.D.; Avina-Bravo, E.G.; Ibarra-Manzano, M.A.; Ruiz-Pinales, J.; Ovalle-Magallanes, E.; Avina-Cervantes, J.G. Retinal Vessel Segmentation Based on a Lightweight U-Net and Reverse Attention. \textit{Mathematics} \textbf{2025}, \textit{13}, 2203. \href{https://doi.org/10.3390/math13132203}{doi:10.3390/math13132203}.



\bibitem{Ho2020}

Ho, J.; Jain, A.; Abbeel, P. Denoising Diffusion Probabilistic Models. In \textit{Advances in Neural Information Processing Systems}; Curran Associates: Red Hook, NY, USA, 2020; Volume 33, pp. 6840--6851.



\bibitem{Dhariwal2021}

Dhariwal, P.; Nichol, A. Diffusion Models Beat GANs on Image Synthesis. In \textit{Advances in Neural Information Processing Systems}; Curran Associates: Red Hook, NY, USA, 2021; Volume 34, pp. 8780--8794.



\bibitem{Rombach2022}

Rombach, R.; Blattmann, A.; Lorenz, D.; Esser, P.; Ommer, B. High-Resolution Image Synthesis with Latent Diffusion Models. In \textit{Proceedings of the IEEE/CVF Conference on Computer Vision and Pattern Recognition (CVPR)}, New Orleans, LA, USA, 18--24 June 2022; pp. 10684--10695. \href{https://doi.org/10.1109/CVPR52688.2022.01042}{doi:10.1109/CVPR52688.2022.01042}.



\bibitem{Soares2006}

Soares, J.V.B.; Leandro, J.J.G.; Cesar, R.M., Jr.; Jelinek, H.F.; Cree, M.J. Retinal Vessel Segmentation Using the 2-D Gabor Wavelet and Supervised Classification. \textit{IEEE Trans. Med. Imaging} \textbf{2006}, \textit{25}, 1214--1222. \href{https://doi.org/10.1109/TMI.2006.879967}{doi:10.1109/TMI.2006.879967}.



\bibitem{Mendonca2006}

Mendon\c{c}a, A.M.; Campilho, A. Segmentation of Retinal Blood Vessels by Combining the Detection of Centerlines and Morphological Reconstruction. \textit{IEEE Trans. Med. Imaging} \textbf{2006}, \textit{25}, 1200--1213. \href{https://doi.org/10.1109/TMI.2006.879955}{doi:10.1109/TMI.2006.879955}.



\bibitem{Ronneberger2015}

Ronneberger, O.; Fischer, P.; Brox, T. U-Net: Convolutional Networks for Biomedical Image Segmentation. In \textit{Medical Image Computing and Computer-Assisted Intervention--MICCAI 2015}; Springer: Cham, Switzerland, 2015; LNCS Volume 9351, pp. 234--241. \href{https://doi.org/10.1007/978-3-319-24574-4_28}{doi:10.1007/978-3-319-24574-4\_28}.



\bibitem{Staal2004}

Staal, J.; Abr\`amoff, M.D.; Niemeijer, M.; Viergever, M.A.; van Ginneken, B. Ridge-Based Vessel Segmentation in Color Images of the Retina. \textit{IEEE Trans. Med. Imaging} \textbf{2004}, \textit{23}, 501--509. \href{https://doi.org/10.1109/TMI.2004.825627}{doi:10.1109/TMI.2004.825627}.



\bibitem{Zhou2018UNetpp}

Zhou, Z.; Siddiquee, M.M.R.; Tajbakhsh, N.; Liang, J. UNet++: A Nested U-Net Architecture for Medical Image Segmentation. In \textit{Deep Learning in Medical Image Analysis and Multimodal Learning for Clinical Decision Support}; Springer: Cham, Switzerland, 2018; LNCS Volume 11045, pp. 3--11. \href{https://doi.org/10.1007/978-3-030-00889-5_1}{doi:10.1007/978-3-030-00889-5\_1}.



\bibitem{Oktay2018AttUNet}

Oktay, O.; Schlemper, J.; Le Folgoc, L.; Lee, M.; Heinrich, M.; Misawa, K.; Mori, K.; McDonagh, S.; Hammerla, N.Y.; Kainz, B.; et al. Attention U-Net: Learning Where to Look for the Pancreas. \textit{arXiv} \textbf{2018}, arXiv:1804.03999. \href{https://doi.org/10.48550/arXiv.1804.03999}{doi:10.48550/arXiv.1804.03999}.



\bibitem{Schlemper2019AG}

Schlemper, J.; Oktay, O.; Schaap, M.; Heinrich, M.; Kainz, B.; Glocker, B.; Rueckert, D. Attention Gated Networks: Learning to Leverage Salient Regions in Medical Images. \textit{Med. Image Anal.} \textbf{2019}, \textit{53}, 197--207. \href{https://doi.org/10.1016/j.media.2019.01.012}{doi:10.1016/j.media.2019.01.012}.



\bibitem{Alom2018R2UNet}

Alom, M.Z.; Yakopcic, C.; Hasan, M.; Taha, T.M.; Asari, V.K. Recurrent Residual U-Net for Medical Image Segmentation. \textit{J. Med. Imaging} \textbf{2019}, \textit{6}, 014006. \href{https://doi.org/10.1117/1.JMI.6.1.014006}{doi:10.1117/1.JMI.6.1.014006}.



\bibitem{Liskowski2016}

Liskowski, P.; Krawiec, K. Segmenting Retinal Blood Vessels with Deep Neural Networks. \textit{IEEE Trans. Med. Imaging} \textbf{2016}, \textit{35}, 2369--2380. \href{https://doi.org/10.1109/TMI.2016.2546227}{doi:10.1109/TMI.2016.2546227}.



\bibitem{Fu2016DeepVessel}

Fu, H.; Xu, Y.; Wong, D.W.K.; Liu, J. Retinal Vessel Segmentation via Deep Learning Network and Fully-Connected Conditional Random Fields. In \textit{Proceedings of the 2016 IEEE 13th International Symposium on Biomedical Imaging (ISBI)}, Prague, Czech Republic, 13--16 April 2016; pp. 698--701. \href{https://doi.org/10.1109/ISBI.2016.7493362}{doi:10.1109/ISBI.2016.7493362}.



\bibitem{Shin2019Connectivity}

Shin, S.Y.; Lee, S.; Yun, I.D.; Lee, K.M. Deep Vessel Segmentation by Learning Graphical Connectivity. \textit{Med. Image Anal.} \textbf{2019}, \textit{58}, 101556. \href{https://doi.org/10.1016/j.media.2019.101556}{doi:10.1016/j.media.2019.101556}.



\bibitem{Song2021}

Song, J.; Meng, C.; Ermon, S. Denoising Diffusion Implicit Models. In \textit{Proceedings of the International Conference on Learning Representations}, Virtual, 3--7 May 2021. Available online: \url{https://openreview.net/forum?id=St1giarCHLP} (accessed on 15 August 2026).



\bibitem{Wolleb2022}

Wolleb, J.; Sandk\"uhler, R.; Bieder, F.; Valmaggia, P.; Cattin, P.C. Diffusion Models for Implicit Image Segmentation Ensembles. In \textit{Proceedings of the 5th International Conference on Medical Imaging with Deep Learning}; PMLR: Cambridge, MA, USA, 2022; Volume 172, pp. 1336--1348.



\bibitem{SegDiff}

Amit, T.; Shaharbany, T.; Nachmani, E.; Wolf, L. SegDiff: Image Segmentation with Diffusion Probabilistic Models. \textit{arXiv} \textbf{2021}, arXiv:2112.00390. \href{https://doi.org/10.48550/arXiv.2112.00390}{doi:10.48550/arXiv.2112.00390}.



\bibitem{MedSegDiff}

Wu, J.; Fu, R.; Fang, H.; Zhang, Y.; Yang, Y.; Xiong, H.; Liu, H.; Xu, Y. MedSegDiff: Medical Image Segmentation with Diffusion Probabilistic Model. In \textit{Medical Imaging with Deep Learning}; PMLR: Cambridge, MA, USA, 2024; Volume 227, pp. 1623--1639.



\bibitem{PDDiffusion}

Guo, X.; Yang, Y.; Ye, C.; Lu, S.; Peng, B.; Huang, H.; Xiang, Y.; Ma, T. Accelerating Diffusion Models via Pre-Segmentation Diffusion Sampling for Medical Image Segmentation. In \textit{Proceedings of the 2023 IEEE 20th International Symposium on Biomedical Imaging (ISBI)}, Cartagena, Colombia, 18--21 April 2023. \href{https://doi.org/10.1109/ISBI53787.2023.10230524}{doi:10.1109/ISBI53787.2023.10230524}.



\bibitem{Kazerouni2023}

Kazerouni, A.; Aghdam, E.K.; Heidari, M.; Azad, R.; Fayyaz, M.; Hacihaliloglu, I.; Merhof, D. Diffusion Models in Medical Imaging: A Comprehensive Survey. \textit{Med. Image Anal.} \textbf{2023}, \textit{88}, 102846. \href{https://doi.org/10.1016/j.media.2023.102846}{doi:10.1016/j.media.2023.102846}.



\bibitem{HiDiffSeg}

Huang, W.; Liu, F. HiDiffSeg: A Hierarchical Diffusion Model for Blood Vessel Segmentation in Retinal Fundus Images. \textit{Expert Syst. Appl.} \textbf{2024}, \textit{253}, 124249. \href{https://doi.org/10.1016/j.eswa.2024.124249}{doi:10.1016/j.eswa.2024.124249}.



\bibitem{RefineCatDiff}

Liu, F.; Yu, H.; Huang, J.; Hua, X.; Liu, H.; Ma, Z.; Wang, Z.; Du, Z.; Sun, L. RefineCatDiff: Toward High-Quality Medical Image Segmentation via a Categorical Diffusion Refinement Framework. \textit{Adv. Intell. Syst.} \textbf{2025}, \textit{7}, e202401125. \href{https://doi.org/10.1002/aisy.202401125}{doi:10.1002/aisy.202401125}.



\bibitem{DiffusionFMD}

Li, G.; Jin, D.; Zheng, Y.; Cui, J.; Gai, W.; Qi, M. A Generic Plug \& Play Diffusion-Based Denosing Module for Medical Image Segmentation. \textit{Neural Netw.} \textbf{2024}, \textit{172}, 106096. \href{https://doi.org/10.1016/j.neunet.2024.106096}{doi:10.1016/j.neunet.2024.106096}.



\bibitem{Hu2018}

Hu, J.; Shen, L.; Sun, G. Squeeze-and-Excitation Networks. In \textit{Proceedings of the IEEE/CVF Conference on Computer Vision and Pattern Recognition (CVPR)}, Salt Lake City, UT, USA, 18--23 June 2018; pp. 7132--7141. \href{https://doi.org/10.1109/CVPR.2018.00745}{doi:10.1109/CVPR.2018.00745}.



\bibitem{Woo2018}

Woo, S.; Park, J.; Lee, J.-Y.; Kweon, I.S. CBAM: Convolutional Block Attention Module. In \textit{Computer Vision--ECCV 2018}; Springer: Cham, Switzerland, 2018; LNCS Volume 11211, pp. 3--19. \href{https://doi.org/10.1007/978-3-030-01234-2_1}{doi:10.1007/978-3-030-01234-2\_1}.



\bibitem{Nichol2021}

Nichol, A.Q.; Dhariwal, P. Improved Denoising Diffusion Probabilistic Models. In \textit{Proceedings of the 38th International Conference on Machine Learning}; PMLR: Cambridge, MA, USA, 2021; Volume 139, pp. 8162--8171.



\bibitem{CHASEDB1}

Fraz, M.M.; Remagnino, P.; Hoppe, A.; Uyyanonvara, B.; Rudnicka, A.R.; Owen, C.G.; Barman, S.A. An Ensemble Classification-Based Approach Applied to Retinal Blood Vessel Segmentation. \textit{IEEE Trans. Biomed. Eng.} \textbf{2012}, \textit{59}, 2538--2548. \href{https://doi.org/10.1109/TBME.2012.2205687}{doi:10.1109/TBME.2012.2205687}. Dataset available online: \url{https://researchinnovation.kingston.ac.uk/en/datasets/chasedb1-retinal-vessel-reference-dataset-4/} (accessed on 15 August 2026).



\bibitem{Odstrcilik2013HRF}

Odstr\v{c}il\'{\i}k, J.; Kola\v{r}, R.; Budai, A.; Hornegger, J.; Jan, J.; Gaz\'arek, J.; Kubena, T.; \v{C}ernosek, P.; Svoboda, O.; Angelopoulou, E. Retinal Vessel Segmentation by Improved Matched Filtering: Evaluation on a New High-Resolution Fundus Image Database. \textit{IET Image Process.} \textbf{2013}, \textit{7}, 373--383. \href{https://doi.org/10.1049/iet-ipr.2012.0455}{doi:10.1049/iet-ipr.2012.0455}. Dataset available online: \url{https://www5.cs.fau.de/research/data/fundus-images/} (accessed on 15 August 2026).

\end{thebibliography}



\ifmdpiclass\PublishersNote{}\fi

%\isPreprints{}{% This command is only used for ``preprints''.



%} % If the paper is ``preprints'', please uncomment this parenthesis.

\end{document}

